
\documentclass{article}
\usepackage{colortbl}
\usepackage{makecell}
\usepackage{multirow}
\usepackage{supertabular}

\begin{document}

\newcounter{utterance}

\centering \large Interaction Transcript for game `clean\_up', experiment `2\_hard\_7obj\_ga', episode 2 with deepseek{-}v4{-}pro.
\vspace{24pt}

{ \footnotesize  \setcounter{utterance}{1}
\setlength{\tabcolsep}{0pt}
\begin{supertabular}{c@{$\;$}|p{.15\linewidth}@{}p{.15\linewidth}p{.15\linewidth}p{.15\linewidth}p{.15\linewidth}p{.15\linewidth}}
    \# & $\;$A & \multicolumn{4}{c}{Game Master} & $\;\:$B\\
    \hline

    \theutterance \stepcounter{utterance}  
    & & \multicolumn{4}{p{0.6\linewidth}}{
        \cellcolor[rgb]{0.9,0.9,0.9}{
            \makecell[{{p{\linewidth}}}]{
                \texttt{\tiny{[P1$\langle$GM]}}
                \texttt{Is mise do mháistir cluiche, agus tá tú ag imirt cluiche comhoibríoch leis an ngreille seo a leanas mar chlár cluiche:} \\
\\ 
\texttt{```} \\
\texttt{ 1234567} \\
\texttt{╔═╤══╤══╗} \\
\texttt{║◌│◌T│◌◌║ 1} \\
\texttt{║◌└┬─┴──╢ 2} \\
\texttt{║◌◌│O◌◌◌║ 3} \\
\texttt{╟─┬┘P◌┌─╢ 4} \\
\texttt{║◌│AD◌│◌║ 5} \\
\texttt{╟─┼─┬─┤M║ 6} \\
\texttt{║◌│S│◌│◌║ 7} \\
\texttt{╚═╧═╧═╧═╝} \\
\\ 
\texttt{```} \\
\\ 
\texttt{* Taispeánann an imeall uachtarach comhordanáidí x ag méadú i dtreo na láimhe deise, agus an t{-}imeall ar dheis comhordanáidí y ag méadú síos.} \\
\texttt{* Cuirtear na rudaí seo a leanas go randamach ar do ghreille: 'P', 'O', 'T', 'S', 'D', 'A', 'M'.} \\
\\ 
\\ 
\texttt{Feiceann an t{-}imreoir eile leagan éagsúil den chlár cluiche, áit a gcuirtear na rudaí in áiteanna randamacha éagsúla. Ní féidir leat clár an imreora eile a fheiceáil, agus ní féidir leo do cheannsa a fheiceáil.} \\
\\ 
\\ 
\texttt{**Sprioc:**} \\
\\ 
\texttt{Caithfidh an dá imreoir na rudaí a bhogadh ar a gcúlra féin ionas go gcríochnóidh rudaí comhionanna ag na comhordanáidí céanna. Caithfidh tú cumarsáid a dhéanamh leis an imreoir eile chun cumraíocht sprice choiteann a aontú.} \\
\\ 
\\ 
\texttt{**Rialacha:**} \\
\\ 
\texttt{* I ngach babhta, is féidir leat ceann amháin go díreach den dá ordú seo a leanas a sheoladh:} \\
\texttt{ 1. `ABAIR: <TEACHTAIREACHT>`: chun teachtaireacht (gach rud suas go dtí an chéad bhriseadh líne eile) a sheoladh chuig an imreoir eile. Cuirfidh mé ar aghaidh í chuig do pháirtí.} \\
\texttt{2. `BOG: <RUD>, (<X>, <Y>)`: chun rud a bhogadh go suíomh nua, áit a bhfuil `<X>` mar an colún agus `<Y>` mar an ró. Cuirfidh mé in iúl duit an raibh do bhogadh rathúil nó nach raibh.} \\
\texttt{* Mura gcloíonn tú leis an bhformáid, nó má sheolann tú roinnt orduithe ag an am céanna, caithfidh mé pionós a ghearradh ort.} \\
\texttt{* Má charnaíonn an dá imreoir níos mó ná 16 pionóis, caillfidh sibh beirt an cluiche.} \\
\texttt{* Tá sé ríthábhachtach go ndéanann tú cumarsáid leis an imreoir eile maidir le do staid sprice! Is é an t{-}aon bhealach *amháin* ar féidir leat do straitéis a chur ar aghaidh chuig an imreoir eile ná an t{-}ordú `ABAIR: <TEACHTAIREACHT>` a úsáid!} \\
\\ 
\\ 
\texttt{**Rudaí a Bhogadh**} \\
\\ 
\texttt{* Ní féidir leat rudaí a bhogadh ach chuig cealla laistigh de theorainneacha na greille. Caithfidh an chill sprice a bheith folamh, is é sin, ní mór di ach an tsiombail '◌' a bheith inti.} \\
\texttt{* Má dhéanann tú iarracht rud a bhogadh go spota nach bhfuil folamh, nó má dhéanann tú iarracht é a bhogadh lasmuigh den ghreille, caithfidh mé pionós a ghearradh ort. Gheobhaidh tú triail eile.} \\
\texttt{* Sula ndéanann tú bogadh, seiceáil faoi dhó go bhfuil an spota sprice folamh, agus nach bhfuil aon litir, fráma ná líne ann!} \\
\\ 
\\ 
\texttt{**Deireadh an Chluiche**} \\
\\ 
\texttt{Má cheapann tú gur shroich tú an sprioc maidir le gach rud a ailíniú, is féidir leat iarraidh ar an imreoir eile an cluiche a chríochnú trí `ABAIR: críochnaithe?` a sheoladh. Má iarrann an t{-}imreoir eile ort an cluiche a chríochnú, agus má fhreagraíonn tú le `ABAIR: críochnaithe!`, tiocfaidh deireadh leis an gcluiche.} \\
\\ 
\texttt{Buaíonn an dá imreoir má chríochnaíonn an cluiche laistigh de 28 babhta, áit a shainmhínítear babhta amháin mar bheirt imreoirí ag seoladh ordú bailí an duine.} \\
\\ 
\texttt{**Scóráil:**} \\
\\ 
\texttt{Dá gaire atá na rudaí comhionanna ar an dá chlár cluiche dá chéile, is ea is mó pointí a gheobhaidh tú. Laghdaíonn pionóis do phointí. An féidir leat an taifead a shárú?} \\
\\ 
\\ 
\texttt{Seol teachtaireacht chuig an imreoir eile chun an cluiche a thosú!} \\
            }
        }
    }
    & \\ \\

    \theutterance \stepcounter{utterance}  
    & \multicolumn{4}{p{0.6\linewidth}}{
        \cellcolor[rgb]{0.9,0.9,0.9}{
            \makecell[{{p{\linewidth}}}]{
                \texttt{\tiny{[P1$\rangle$GM]}}
                \texttt{ABAIR: Dia duit! Feicfidh muid an féidir linn ár rudaí a ailíniú. Seo iad mo rudaí agus a suímh: P ag (3,4), O ag (3,3), T ag (3,1), S ag (2,7), D ag (3,5), A ag (2,5), M ag (7,6). Cad iad do shuímh{-}se?} \\
            }
        }
    }
    & & \\ \\

    \theutterance \stepcounter{utterance}  
    & & & \multicolumn{2}{p{0.3\linewidth}}{
        \cellcolor[rgb]{0.9,0.9,0.9}{
            \makecell[{{p{\linewidth}}}]{
                \texttt{\tiny{[GM$|$GM]}}
                \texttt{ABAIR: Dia duit! Feicfidh muid an féidir linn ár rudaí a ailíniú. Seo iad mo rudaí agus a suímh: P ag (3,4), O ag (3,3), T ag (3,1), S ag (2,7), D ag (3,5), A ag (2,5), M ag (7,6). Cad iad do shuímh{-}se?} \\
            }
        }
    }
    & & \\ \\

    \theutterance \stepcounter{utterance}  
    & & \multicolumn{4}{p{0.6\linewidth}}{
        \cellcolor[rgb]{0.9,0.9,0.9}{
            \makecell[{{p{\linewidth}}}]{
                \texttt{\tiny{[GM$\rangle$P2]}}
                \texttt{Is mise do mháistir cluiche, agus tá tú ag imirt cluiche comhoibríoch leis an ngreille seo a leanas mar chlár cluiche:} \\
\\ 
\texttt{```} \\
\texttt{ 1234567} \\
\texttt{╔═╤══╤══╗} \\
\texttt{║◌│O◌│◌◌║ 1} \\
\texttt{║◌└┬─┴──╢ 2} \\
\texttt{║◌P│◌D◌◌║ 3} \\
\texttt{╟─┬┘◌◌┌─╢ 4} \\
\texttt{║◌│◌◌◌│A║ 5} \\
\texttt{╟─┼─┬─┤S║ 6} \\
\texttt{║M│◌│◌│T║ 7} \\
\texttt{╚═╧═╧═╧═╝} \\
\\ 
\texttt{```} \\
\\ 
\texttt{* Taispeánann an imeall uachtarach comhordanáidí x ag méadú i dtreo na láimhe deise, agus an t{-}imeall ar dheis comhordanáidí y ag méadú síos.} \\
\texttt{* Cuirtear na rudaí seo a leanas go randamach ar do ghreille: 'P', 'O', 'T', 'S', 'D', 'A', 'M'.} \\
\\ 
\\ 
\texttt{Feiceann an t{-}imreoir eile leagan éagsúil den chlár cluiche, áit a gcuirtear na rudaí in áiteanna randamacha éagsúla. Ní féidir leat clár an imreora eile a fheiceáil, agus ní féidir leo do cheannsa a fheiceáil.} \\
\\ 
\\ 
\texttt{**Sprioc:**} \\
\\ 
\texttt{Caithfidh an dá imreoir na rudaí a bhogadh ar a gcúlra féin ionas go gcríochnóidh rudaí comhionanna ag na comhordanáidí céanna. Caithfidh tú cumarsáid a dhéanamh leis an imreoir eile chun cumraíocht sprice choiteann a aontú.} \\
\\ 
\\ 
\texttt{**Rialacha:**} \\
\\ 
\texttt{* I ngach babhta, is féidir leat ceann amháin go díreach den dá ordú seo a leanas a sheoladh:} \\
\texttt{ 1. `ABAIR: <TEACHTAIREACHT>`: chun teachtaireacht (gach rud suas go dtí an chéad bhriseadh líne eile) a sheoladh chuig an imreoir eile. Cuirfidh mé ar aghaidh í chuig do pháirtí.} \\
\texttt{2. `BOG: <RUD>, (<X>, <Y>)`: chun rud a bhogadh go suíomh nua, áit a bhfuil `<X>` mar an colún agus `<Y>` mar an ró. Cuirfidh mé in iúl duit an raibh do bhogadh rathúil nó nach raibh.} \\
\texttt{* Mura gcloíonn tú leis an bhformáid, nó má sheolann tú roinnt orduithe ag an am céanna, caithfidh mé pionós a ghearradh ort.} \\
\texttt{* Má charnaíonn an dá imreoir níos mó ná 16 pionóis, caillfidh sibh beirt an cluiche.} \\
\texttt{* Tá sé ríthábhachtach go ndéanann tú cumarsáid leis an imreoir eile maidir le do staid sprice! Is é an t{-}aon bhealach *amháin* ar féidir leat do straitéis a chur ar aghaidh chuig an imreoir eile ná an t{-}ordú `ABAIR: <TEACHTAIREACHT>` a úsáid!} \\
\\ 
\\ 
\texttt{**Rudaí a Bhogadh**} \\
\\ 
\texttt{* Ní féidir leat rudaí a bhogadh ach chuig cealla laistigh de theorainneacha na greille. Caithfidh an chill sprice a bheith folamh, is é sin, ní mór di ach an tsiombail '◌' a bheith inti.} \\
\texttt{* Má dhéanann tú iarracht rud a bhogadh go spota nach bhfuil folamh, nó má dhéanann tú iarracht é a bhogadh lasmuigh den ghreille, caithfidh mé pionós a ghearradh ort. Gheobhaidh tú triail eile.} \\
\texttt{* Sula ndéanann tú bogadh, seiceáil faoi dhó go bhfuil an spota sprice folamh, agus nach bhfuil aon litir, fráma ná líne ann!} \\
\\ 
\\ 
\texttt{**Deireadh an Chluiche**} \\
\\ 
\texttt{Má cheapann tú gur shroich tú an sprioc maidir le gach rud a ailíniú, is féidir leat iarraidh ar an imreoir eile an cluiche a chríochnú trí `ABAIR: críochnaithe?` a sheoladh. Má iarrann an t{-}imreoir eile ort an cluiche a chríochnú, agus má fhreagraíonn tú le `ABAIR: críochnaithe!`, tiocfaidh deireadh leis an gcluiche.} \\
\\ 
\texttt{Buaíonn an dá imreoir má chríochnaíonn an cluiche laistigh de 28 babhta, áit a shainmhínítear babhta amháin mar bheirt imreoirí ag seoladh ordú bailí an duine.} \\
\\ 
\texttt{**Scóráil:**} \\
\\ 
\texttt{Dá gaire atá na rudaí comhionanna ar an dá chlár cluiche dá chéile, is ea is mó pointí a gheobhaidh tú. Laghdaíonn pionóis do phointí. An féidir leat an taifead a shárú?} \\
\\ 
\texttt{Thosaigh an t{-}imreoir eile an cluiche trí an teachtaireacht seo a sheoladh:} \\
\\ 
\texttt{"Dia duit! Feicfidh muid an féidir linn ár rudaí a ailíniú. Seo iad mo rudaí agus a suímh: P ag (3,4), O ag (3,3), T ag (3,1), S ag (2,7), D ag (3,5), A ag (2,5), M ag (7,6). Cad iad do shuímh{-}se?"} \\
\\ 
\texttt{Cad é do chéad ordú?} \\
            }
        }
    }
    & \\ \\

    \theutterance \stepcounter{utterance}  
    & & & \multicolumn{4}{p{0.6\linewidth}}{
        \cellcolor[rgb]{0.9,0.9,0.9}{
            \makecell[{{p{\linewidth}}}]{
                \texttt{\tiny{[GM$\langle$P2]}}
                \texttt{ABAIR: Dia duit! Seo iad mo shuímh: P ag (2,3), O ag (2,1), T ag (7,7), S ag (6,6), D ag (5,3), A ag (7,5), M ag (1,7). Lig dúinn na suímh a chur i gcomparáid agus teacht ar chomhaontú.} \\
            }
        }
    }
     \\ \\

    \theutterance \stepcounter{utterance}  
    & & & \multicolumn{2}{p{0.3\linewidth}}{
        \cellcolor[rgb]{0.9,0.9,0.9}{
            \makecell[{{p{\linewidth}}}]{
                \texttt{\tiny{[GM$|$GM]}}
                \texttt{ABAIR: Dia duit! Seo iad mo shuímh: P ag (2,3), O ag (2,1), T ag (7,7), S ag (6,6), D ag (5,3), A ag (7,5), M ag (1,7). Lig dúinn na suímh a chur i gcomparáid agus teacht ar chomhaontú.} \\
            }
        }
    }
    & & \\ \\

    \theutterance \stepcounter{utterance}  
    & & \multicolumn{4}{p{0.6\linewidth}}{
        \cellcolor[rgb]{0.9,0.9,0.9}{
            \makecell[{{p{\linewidth}}}]{
                \texttt{\tiny{[P1$\langle$GM]}}
                \texttt{Cuireadh do theachtaireacht ar aghaidh chuig an imreoir eile.} \\
\\ 
\\ 
\texttt{Tá sibh ag imirt babhta 1 faoi láthair as uasmhéid 28.} \\
\texttt{Tá 0 as 16 pionós carntha agaibh le chéile.} \\
\texttt{Sheol an t{-}imreoir eile an teachtaireacht seo:} \\
\\ 
\\ 
\texttt{"Dia duit! Seo iad mo shuímh: P ag (2,3), O ag (2,1), T ag (7,7), S ag (6,6), D ag (5,3), A ag (7,5), M ag (1,7). Lig dúinn na suímh a chur i gcomparáid agus teacht ar chomhaontú."} \\
\\ 
\\ 
\texttt{Cad é do chéad ordú eile?} \\
            }
        }
    }
    & \\ \\

    \theutterance \stepcounter{utterance}  
    & \multicolumn{4}{p{0.6\linewidth}}{
        \cellcolor[rgb]{0.9,0.9,0.9}{
            \makecell[{{p{\linewidth}}}]{
                \texttt{\tiny{[P1$\rangle$GM]}}
                \texttt{ABAIR: Go raibh maith agat! Déanaimis comparáid:} \\
\\ 
\texttt{Mo shuímh: P(3,4), O(3,3), T(3,1), S(2,7), D(3,5), A(2,5), M(7,6)} \\
\texttt{Do shuímh: P(2,3), O(2,1), T(7,7), S(6,6), D(5,3), A(7,5), M(1,7)} \\
\\ 
\texttt{Caithfimid spriocanna comónta a roghnú. Molaim an lár mar sprioc: P(4,4), O(4,3), T(4,2), S(3,4), D(5,4), A(4,5), M(4,6) {-} is féidir linn patrún ingearach a dhéanamh sa cholún 4. Nó an bhfuil smaoineamh eile agat? Tá na rudaí scaipthe go leor againn beirt.} \\
            }
        }
    }
    & & \\ \\

    \theutterance \stepcounter{utterance}  
    & & & \multicolumn{2}{p{0.3\linewidth}}{
        \cellcolor[rgb]{0.9,0.9,0.9}{
            \makecell[{{p{\linewidth}}}]{
                \texttt{\tiny{[GM$|$GM]}}
                \texttt{ABAIR: Go raibh maith agat! Déanaimis comparáid:} \\
\\ 
\texttt{Mo shuímh: P(3,4), O(3,3), T(3,1), S(2,7), D(3,5), A(2,5), M(7,6)} \\
\texttt{Do shuímh: P(2,3), O(2,1), T(7,7), S(6,6), D(5,3), A(7,5), M(1,7)} \\
\\ 
\texttt{Caithfimid spriocanna comónta a roghnú. Molaim an lár mar sprioc: P(4,4), O(4,3), T(4,2), S(3,4), D(5,4), A(4,5), M(4,6) {-} is féidir linn patrún ingearach a dhéanamh sa cholún 4. Nó an bhfuil smaoineamh eile agat? Tá na rudaí scaipthe go leor againn beirt.} \\
            }
        }
    }
    & & \\ \\

    \theutterance \stepcounter{utterance}  
    & & & \multicolumn{2}{p{0.3\linewidth}}{
        \cellcolor[rgb]{0.9,0.9,0.9}{
            \makecell[{{p{\linewidth}}}]{
                \texttt{\tiny{[GM$|$GM]}}
                \texttt{Invalid format: tail is not empty: '} \\
\\ 
\texttt{Mo shuímh: P(3,4), O(3,3), T(3,1), S(2,7), D(3,5), A(2,5), M(7,6)} \\
\texttt{Do shuímh: P(2,3), O(2,1), T(7,7), S(6,6), D(5,3), A(7,5), M(1,7)} \\
\\ 
\texttt{Caithfimid spriocanna comónta a roghnú. Molaim an lár mar sprioc: P(4,4), O(4,3), T(4,2), S(3,4), D(5,4), A(4,5), M(4,6) {-} is féidir linn patrún ingearach a dhéanamh sa cholún 4. Nó an bhfuil smaoineamh eile agat? Tá na rudaí scaipthe go leor againn beirt.} \\
            }
        }
    }
    & & \\ \\

    \theutterance \stepcounter{utterance}  
    & & \multicolumn{4}{p{0.6\linewidth}}{
        \cellcolor[rgb]{0.9,0.9,0.9}{
            \makecell[{{p{\linewidth}}}]{
                \texttt{\tiny{[P1$\langle$GM]}}
                \texttt{Pionós: Ní cheadaítear aon rud a bheith i do theachtaireacht tar éis an ordaithe!} \\
\\ 
\texttt{Cinntigh nach bhfuil i do fhreagra ach `ABAIR: <TEACHTAIREACHT>` nó `BOG: <RUD>, (<X>, <Y>)`, agus tada eile!} \\
\\ 
\texttt{Tá 1 as 16 pionós carntha agaibh le chéile.} \\
\texttt{Déan iarracht arís, le do thoil!} \\
            }
        }
    }
    & \\ \\

    \theutterance \stepcounter{utterance}  
    & \multicolumn{4}{p{0.6\linewidth}}{
        \cellcolor[rgb]{0.9,0.9,0.9}{
            \makecell[{{p{\linewidth}}}]{
                \texttt{\tiny{[P1$\rangle$GM]}}
                \texttt{ABAIR: Go raibh maith agat! Déanaimis comparáid: Mo shuímh: P(3,4), O(3,3), T(3,1), S(2,7), D(3,5), A(2,5), M(7,6). Do shuímh: P(2,3), O(2,1), T(7,7), S(6,6), D(5,3), A(7,5), M(1,7). An féidir linn teacht ar chomhaontú faoi sprioc{-}shuímh? Mar shampla, an bhfuil ciall le P(4,4), O(4,3), T(4,2), S(3,4), D(5,4), A(4,5), M(4,6)? Nó an bhfuil moladh eile agat?} \\
            }
        }
    }
    & & \\ \\

    \theutterance \stepcounter{utterance}  
    & & & \multicolumn{2}{p{0.3\linewidth}}{
        \cellcolor[rgb]{0.9,0.9,0.9}{
            \makecell[{{p{\linewidth}}}]{
                \texttt{\tiny{[GM$|$GM]}}
                \texttt{ABAIR: Go raibh maith agat! Déanaimis comparáid: Mo shuímh: P(3,4), O(3,3), T(3,1), S(2,7), D(3,5), A(2,5), M(7,6). Do shuímh: P(2,3), O(2,1), T(7,7), S(6,6), D(5,3), A(7,5), M(1,7). An féidir linn teacht ar chomhaontú faoi sprioc{-}shuímh? Mar shampla, an bhfuil ciall le P(4,4), O(4,3), T(4,2), S(3,4), D(5,4), A(4,5), M(4,6)? Nó an bhfuil moladh eile agat?} \\
            }
        }
    }
    & & \\ \\

    \theutterance \stepcounter{utterance}  
    & & \multicolumn{4}{p{0.6\linewidth}}{
        \cellcolor[rgb]{0.9,0.9,0.9}{
            \makecell[{{p{\linewidth}}}]{
                \texttt{\tiny{[GM$\rangle$P2]}}
                \texttt{Cuireadh do theachtaireacht ar aghaidh chuig an imreoir eile.} \\
\\ 
\\ 
\texttt{Tá sibh ag imirt babhta 2 faoi láthair as uasmhéid 28.} \\
\texttt{Tá 1 as 16 pionós carntha agaibh le chéile.} \\
\texttt{Sheol an t{-}imreoir eile an teachtaireacht seo:} \\
\\ 
\\ 
\texttt{"Go raibh maith agat! Déanaimis comparáid: Mo shuímh: P(3,4), O(3,3), T(3,1), S(2,7), D(3,5), A(2,5), M(7,6). Do shuímh: P(2,3), O(2,1), T(7,7), S(6,6), D(5,3), A(7,5), M(1,7). An féidir linn teacht ar chomhaontú faoi sprioc{-}shuímh? Mar shampla, an bhfuil ciall le P(4,4), O(4,3), T(4,2), S(3,4), D(5,4), A(4,5), M(4,6)? Nó an bhfuil moladh eile agat?"} \\
\\ 
\\ 
\texttt{Cad é do chéad ordú eile?} \\
            }
        }
    }
    & \\ \\

    \theutterance \stepcounter{utterance}  
    & & & \multicolumn{4}{p{0.6\linewidth}}{
        \cellcolor[rgb]{0.9,0.9,0.9}{
            \makecell[{{p{\linewidth}}}]{
                \texttt{\tiny{[GM$\langle$P2]}}
                \texttt{ABAIR: Is maith liom do smaoineamh sprioc{-}shuímh a roghnú! Ach déanaimis seiceáil ar dtús an bhfuil na spriocanna sin folamh ar ár gcláir. Ar do chlár, an bhfuil na cealla seo folamh: (4,4), (4,3), (4,2), (3,4), (5,4), (4,5), (4,6)? Ar mo chlár féin, tá (4,4) folamh, (4,3) folamh, (4,2) folamh, (3,4) tá líne ann, (5,4) folamh, (4,5) folamh, (4,6) folamh. B'fhéidir go gcaithfimid na sprioc{-}shuímh a choigeartú beagán mar gheall ar na constaicí. An bhféadfá do chlár a sheiceáil?} \\
            }
        }
    }
     \\ \\

    \theutterance \stepcounter{utterance}  
    & & & \multicolumn{2}{p{0.3\linewidth}}{
        \cellcolor[rgb]{0.9,0.9,0.9}{
            \makecell[{{p{\linewidth}}}]{
                \texttt{\tiny{[GM$|$GM]}}
                \texttt{ABAIR: Is maith liom do smaoineamh sprioc{-}shuímh a roghnú! Ach déanaimis seiceáil ar dtús an bhfuil na spriocanna sin folamh ar ár gcláir. Ar do chlár, an bhfuil na cealla seo folamh: (4,4), (4,3), (4,2), (3,4), (5,4), (4,5), (4,6)? Ar mo chlár féin, tá (4,4) folamh, (4,3) folamh, (4,2) folamh, (3,4) tá líne ann, (5,4) folamh, (4,5) folamh, (4,6) folamh. B'fhéidir go gcaithfimid na sprioc{-}shuímh a choigeartú beagán mar gheall ar na constaicí. An bhféadfá do chlár a sheiceáil?} \\
            }
        }
    }
    & & \\ \\

    \theutterance \stepcounter{utterance}  
    & & \multicolumn{4}{p{0.6\linewidth}}{
        \cellcolor[rgb]{0.9,0.9,0.9}{
            \makecell[{{p{\linewidth}}}]{
                \texttt{\tiny{[P1$\langle$GM]}}
                \texttt{Cuireadh do theachtaireacht ar aghaidh chuig an imreoir eile.} \\
\\ 
\\ 
\texttt{Tá sibh ag imirt babhta 2 faoi láthair as uasmhéid 28.} \\
\texttt{Tá 1 as 16 pionós carntha agaibh le chéile.} \\
\texttt{Sheol an t{-}imreoir eile an teachtaireacht seo:} \\
\\ 
\\ 
\texttt{"Is maith liom do smaoineamh sprioc{-}shuímh a roghnú! Ach déanaimis seiceáil ar dtús an bhfuil na spriocanna sin folamh ar ár gcláir. Ar do chlár, an bhfuil na cealla seo folamh: (4,4), (4,3), (4,2), (3,4), (5,4), (4,5), (4,6)? Ar mo chlár féin, tá (4,4) folamh, (4,3) folamh, (4,2) folamh, (3,4) tá líne ann, (5,4) folamh, (4,5) folamh, (4,6) folamh. B'fhéidir go gcaithfimid na sprioc{-}shuímh a choigeartú beagán mar gheall ar na constaicí. An bhféadfá do chlár a sheiceáil?"} \\
\\ 
\\ 
\texttt{Cad é do chéad ordú eile?} \\
            }
        }
    }
    & \\ \\

    \theutterance \stepcounter{utterance}  
    & \multicolumn{4}{p{0.6\linewidth}}{
        \cellcolor[rgb]{0.9,0.9,0.9}{
            \makecell[{{p{\linewidth}}}]{
                \texttt{\tiny{[P1$\rangle$GM]}}
                \texttt{ABAIR: Rinne mé seiceáil ar mo chlár. Tá (4,4) folamh, (4,3) tá líne ingearach ann, (4,2) tá líne ingearach ann, (3,4) folamh, (5,4) folamh, (4,5) folamh, (4,6) folamh. Mar sin tá constaicí againn beirt ag (4,3) agus (4,2) domsa, agus ag (3,4) duitse. An féidir linn na sprioc{-}shuímh a bhogadh? Mar shampla, bain úsáid as colún 5 mar mhalairt: P(5,4), O(5,3), T(5,2), S(4,4), D(6,4), A(5,5), M(5,6)? Nó an bhfuil moladh níos fearr agat?} \\
            }
        }
    }
    & & \\ \\

    \theutterance \stepcounter{utterance}  
    & & & \multicolumn{2}{p{0.3\linewidth}}{
        \cellcolor[rgb]{0.9,0.9,0.9}{
            \makecell[{{p{\linewidth}}}]{
                \texttt{\tiny{[GM$|$GM]}}
                \texttt{ABAIR: Rinne mé seiceáil ar mo chlár. Tá (4,4) folamh, (4,3) tá líne ingearach ann, (4,2) tá líne ingearach ann, (3,4) folamh, (5,4) folamh, (4,5) folamh, (4,6) folamh. Mar sin tá constaicí againn beirt ag (4,3) agus (4,2) domsa, agus ag (3,4) duitse. An féidir linn na sprioc{-}shuímh a bhogadh? Mar shampla, bain úsáid as colún 5 mar mhalairt: P(5,4), O(5,3), T(5,2), S(4,4), D(6,4), A(5,5), M(5,6)? Nó an bhfuil moladh níos fearr agat?} \\
            }
        }
    }
    & & \\ \\

    \theutterance \stepcounter{utterance}  
    & & \multicolumn{4}{p{0.6\linewidth}}{
        \cellcolor[rgb]{0.9,0.9,0.9}{
            \makecell[{{p{\linewidth}}}]{
                \texttt{\tiny{[GM$\rangle$P2]}}
                \texttt{Cuireadh do theachtaireacht ar aghaidh chuig an imreoir eile.} \\
\\ 
\\ 
\texttt{Tá sibh ag imirt babhta 3 faoi láthair as uasmhéid 28.} \\
\texttt{Tá 1 as 16 pionós carntha agaibh le chéile.} \\
\texttt{Sheol an t{-}imreoir eile an teachtaireacht seo:} \\
\\ 
\\ 
\texttt{"Rinne mé seiceáil ar mo chlár. Tá (4,4) folamh, (4,3) tá líne ingearach ann, (4,2) tá líne ingearach ann, (3,4) folamh, (5,4) folamh, (4,5) folamh, (4,6) folamh. Mar sin tá constaicí againn beirt ag (4,3) agus (4,2) domsa, agus ag (3,4) duitse. An féidir linn na sprioc{-}shuímh a bhogadh? Mar shampla, bain úsáid as colún 5 mar mhalairt: P(5,4), O(5,3), T(5,2), S(4,4), D(6,4), A(5,5), M(5,6)? Nó an bhfuil moladh níos fearr agat?"} \\
\\ 
\\ 
\texttt{Cad é do chéad ordú eile?} \\
            }
        }
    }
    & \\ \\

    \theutterance \stepcounter{utterance}  
    & & & \multicolumn{4}{p{0.6\linewidth}}{
        \cellcolor[rgb]{0.9,0.9,0.9}{
            \makecell[{{p{\linewidth}}}]{
                \texttt{\tiny{[GM$\langle$P2]}}
                \texttt{ABAIR: Seiceáil mé mo chlár do na suímh nua sin: (5,4) folamh, (5,3) tá D ann faoi láthair ach is féidir é a bhogadh, (5,2) folamh, (4,4) folamh, (6,4) folamh, (5,5) folamh, (5,6) folamh. Is féidir linn triail a bhaint as an bplean seo! An féidir leat tosú ag bogadh do rudaí chuig na sprioc{-}shuímh nua? Cuirfidh mé tús le mo chuid bogtha freisin. Mar shampla, ba chóir dom P a bhogadh ó (2,3) go (5,4). An bhfuil tú réidh le tosú?} \\
            }
        }
    }
     \\ \\

    \theutterance \stepcounter{utterance}  
    & & & \multicolumn{2}{p{0.3\linewidth}}{
        \cellcolor[rgb]{0.9,0.9,0.9}{
            \makecell[{{p{\linewidth}}}]{
                \texttt{\tiny{[GM$|$GM]}}
                \texttt{ABAIR: Seiceáil mé mo chlár do na suímh nua sin: (5,4) folamh, (5,3) tá D ann faoi láthair ach is féidir é a bhogadh, (5,2) folamh, (4,4) folamh, (6,4) folamh, (5,5) folamh, (5,6) folamh. Is féidir linn triail a bhaint as an bplean seo! An féidir leat tosú ag bogadh do rudaí chuig na sprioc{-}shuímh nua? Cuirfidh mé tús le mo chuid bogtha freisin. Mar shampla, ba chóir dom P a bhogadh ó (2,3) go (5,4). An bhfuil tú réidh le tosú?} \\
            }
        }
    }
    & & \\ \\

    \theutterance \stepcounter{utterance}  
    & & \multicolumn{4}{p{0.6\linewidth}}{
        \cellcolor[rgb]{0.9,0.9,0.9}{
            \makecell[{{p{\linewidth}}}]{
                \texttt{\tiny{[P1$\langle$GM]}}
                \texttt{Cuireadh do theachtaireacht ar aghaidh chuig an imreoir eile.} \\
\\ 
\\ 
\texttt{Tá sibh ag imirt babhta 3 faoi láthair as uasmhéid 28.} \\
\texttt{Tá 1 as 16 pionós carntha agaibh le chéile.} \\
\texttt{Sheol an t{-}imreoir eile an teachtaireacht seo:} \\
\\ 
\\ 
\texttt{"Seiceáil mé mo chlár do na suímh nua sin: (5,4) folamh, (5,3) tá D ann faoi láthair ach is féidir é a bhogadh, (5,2) folamh, (4,4) folamh, (6,4) folamh, (5,5) folamh, (5,6) folamh. Is féidir linn triail a bhaint as an bplean seo! An féidir leat tosú ag bogadh do rudaí chuig na sprioc{-}shuímh nua? Cuirfidh mé tús le mo chuid bogtha freisin. Mar shampla, ba chóir dom P a bhogadh ó (2,3) go (5,4). An bhfuil tú réidh le tosú?"} \\
\\ 
\\ 
\texttt{Cad é do chéad ordú eile?} \\
            }
        }
    }
    & \\ \\

    \theutterance \stepcounter{utterance}  
    & \multicolumn{4}{p{0.6\linewidth}}{
        \cellcolor[rgb]{0.9,0.9,0.9}{
            \makecell[{{p{\linewidth}}}]{
                \texttt{\tiny{[P1$\rangle$GM]}}
                \texttt{BOG: P, (5,4)} \\
            }
        }
    }
    & & \\ \\

    \theutterance \stepcounter{utterance}  
    & & & \multicolumn{2}{p{0.3\linewidth}}{
        \cellcolor[rgb]{0.9,0.9,0.9}{
            \makecell[{{p{\linewidth}}}]{
                \texttt{\tiny{[GM$|$GM]}}
                \texttt{BOG: P, (5,4)} \\
            }
        }
    }
    & & \\ \\

    \theutterance \stepcounter{utterance}  
    & & & \multicolumn{2}{p{0.3\linewidth}}{
        \cellcolor[rgb]{0.9,0.9,0.9}{
            \makecell[{{p{\linewidth}}}]{
                \texttt{\tiny{[GM$|$GM]}}
                \texttt{P moved to (5, 4)} \\
            }
        }
    }
    & & \\ \\

    \theutterance \stepcounter{utterance}  
    & & \multicolumn{4}{p{0.6\linewidth}}{
        \cellcolor[rgb]{0.9,0.9,0.9}{
            \makecell[{{p{\linewidth}}}]{
                \texttt{\tiny{[GM$\rangle$P2]}}
                \texttt{Cuireadh do theachtaireacht ar aghaidh chuig an imreoir eile.} \\
\\ 
\\ 
\texttt{Tá sibh ag imirt babhta 4 faoi láthair as uasmhéid 28.} \\
\texttt{Tá 1 as 16 pionós carntha agaibh le chéile.} \\
\texttt{Bhog an t{-}imreoir eile réad ar a ngreille.} \\
\\ 
\\ 
\texttt{Cad é do chéad ordú eile?} \\
            }
        }
    }
    & \\ \\

    \theutterance \stepcounter{utterance}  
    & & & \multicolumn{4}{p{0.6\linewidth}}{
        \cellcolor[rgb]{0.9,0.9,0.9}{
            \makecell[{{p{\linewidth}}}]{
                \texttt{\tiny{[GM$\langle$P2]}}
                \texttt{BOG: P, (5,4)} \\
            }
        }
    }
     \\ \\

    \theutterance \stepcounter{utterance}  
    & & & \multicolumn{2}{p{0.3\linewidth}}{
        \cellcolor[rgb]{0.9,0.9,0.9}{
            \makecell[{{p{\linewidth}}}]{
                \texttt{\tiny{[GM$|$GM]}}
                \texttt{BOG: P, (5,4)} \\
            }
        }
    }
    & & \\ \\

    \theutterance \stepcounter{utterance}  
    & & & \multicolumn{2}{p{0.3\linewidth}}{
        \cellcolor[rgb]{0.9,0.9,0.9}{
            \makecell[{{p{\linewidth}}}]{
                \texttt{\tiny{[GM$|$GM]}}
                \texttt{P moved to (5, 4)} \\
            }
        }
    }
    & & \\ \\

    \theutterance \stepcounter{utterance}  
    & & \multicolumn{4}{p{0.6\linewidth}}{
        \cellcolor[rgb]{0.9,0.9,0.9}{
            \makecell[{{p{\linewidth}}}]{
                \texttt{\tiny{[P1$\langle$GM]}}
                \texttt{Bogadh 'P' go (5,4) go rathúil. Seo mar atá do ghreille nuashonraithe:} \\
\texttt{```} \\
\texttt{ 1234567} \\
\texttt{╔═╤══╤══╗} \\
\texttt{║◌│◌T│◌◌║ 1} \\
\texttt{║◌└┬─┴──╢ 2} \\
\texttt{║◌◌│O◌◌◌║ 3} \\
\texttt{╟─┬┘◌P┌─╢ 4} \\
\texttt{║◌│AD◌│◌║ 5} \\
\texttt{╟─┼─┬─┤M║ 6} \\
\texttt{║◌│S│◌│◌║ 7} \\
\texttt{╚═╧═╧═╧═╝} \\
\texttt{```} \\
\texttt{Tá sibh ag imirt babhta 4 faoi láthair as uasmhéid 28.} \\
\texttt{Tá 1 as 16 pionós carntha agaibh le chéile.} \\
\texttt{Bhog an t{-}imreoir eile réad ar a ngreille.} \\
\\ 
\\ 
\texttt{Cad é do chéad ordú eile?} \\
            }
        }
    }
    & \\ \\

    \theutterance \stepcounter{utterance}  
    & \multicolumn{4}{p{0.6\linewidth}}{
        \cellcolor[rgb]{0.9,0.9,0.9}{
            \makecell[{{p{\linewidth}}}]{
                \texttt{\tiny{[P1$\rangle$GM]}}
                \texttt{BOG: O, (5,3)} \\
            }
        }
    }
    & & \\ \\

    \theutterance \stepcounter{utterance}  
    & & & \multicolumn{2}{p{0.3\linewidth}}{
        \cellcolor[rgb]{0.9,0.9,0.9}{
            \makecell[{{p{\linewidth}}}]{
                \texttt{\tiny{[GM$|$GM]}}
                \texttt{BOG: O, (5,3)} \\
            }
        }
    }
    & & \\ \\

    \theutterance \stepcounter{utterance}  
    & & & \multicolumn{2}{p{0.3\linewidth}}{
        \cellcolor[rgb]{0.9,0.9,0.9}{
            \makecell[{{p{\linewidth}}}]{
                \texttt{\tiny{[GM$|$GM]}}
                \texttt{O moved to (5, 3)} \\
            }
        }
    }
    & & \\ \\

    \theutterance \stepcounter{utterance}  
    & & \multicolumn{4}{p{0.6\linewidth}}{
        \cellcolor[rgb]{0.9,0.9,0.9}{
            \makecell[{{p{\linewidth}}}]{
                \texttt{\tiny{[GM$\rangle$P2]}}
                \texttt{Bogadh 'P' go (5,4) go rathúil. Seo mar atá do ghreille nuashonraithe:} \\
\texttt{```} \\
\texttt{ 1234567} \\
\texttt{╔═╤══╤══╗} \\
\texttt{║◌│O◌│◌◌║ 1} \\
\texttt{║◌└┬─┴──╢ 2} \\
\texttt{║◌◌│◌D◌◌║ 3} \\
\texttt{╟─┬┘◌P┌─╢ 4} \\
\texttt{║◌│◌◌◌│A║ 5} \\
\texttt{╟─┼─┬─┤S║ 6} \\
\texttt{║M│◌│◌│T║ 7} \\
\texttt{╚═╧═╧═╧═╝} \\
\texttt{```} \\
\texttt{Tá sibh ag imirt babhta 5 faoi láthair as uasmhéid 28.} \\
\texttt{Tá 1 as 16 pionós carntha agaibh le chéile.} \\
\texttt{Bhog an t{-}imreoir eile réad ar a ngreille.} \\
\\ 
\\ 
\texttt{Cad é do chéad ordú eile?} \\
            }
        }
    }
    & \\ \\

    \theutterance \stepcounter{utterance}  
    & & & \multicolumn{4}{p{0.6\linewidth}}{
        \cellcolor[rgb]{0.9,0.9,0.9}{
            \makecell[{{p{\linewidth}}}]{
                \texttt{\tiny{[GM$\langle$P2]}}
                \texttt{BOG: O, (5,3)} \\
            }
        }
    }
     \\ \\

    \theutterance \stepcounter{utterance}  
    & & & \multicolumn{2}{p{0.3\linewidth}}{
        \cellcolor[rgb]{0.9,0.9,0.9}{
            \makecell[{{p{\linewidth}}}]{
                \texttt{\tiny{[GM$|$GM]}}
                \texttt{BOG: O, (5,3)} \\
            }
        }
    }
    & & \\ \\

    \theutterance \stepcounter{utterance}  
    & & & \multicolumn{2}{p{0.3\linewidth}}{
        \cellcolor[rgb]{0.9,0.9,0.9}{
            \makecell[{{p{\linewidth}}}]{
                \texttt{\tiny{[GM$|$GM]}}
                \texttt{Pionós: níl (5,3) folamh, ach tá 'D' ann.} \\
\\ 
\texttt{Tá 2 as 16 pionós carntha agaibh le chéile.} \\
\texttt{Déan iarracht arís, le do thoil!} \\
            }
        }
    }
    & & \\ \\

    \theutterance \stepcounter{utterance}  
    & & \multicolumn{4}{p{0.6\linewidth}}{
        \cellcolor[rgb]{0.9,0.9,0.9}{
            \makecell[{{p{\linewidth}}}]{
                \texttt{\tiny{[GM$\rangle$P2]}}
                \texttt{Pionós: níl (5,3) folamh, ach tá 'D' ann.} \\
\\ 
\texttt{Tá 2 as 16 pionós carntha agaibh le chéile.} \\
\texttt{Déan iarracht arís, le do thoil!} \\
            }
        }
    }
    & \\ \\

    \theutterance \stepcounter{utterance}  
    & & & \multicolumn{4}{p{0.6\linewidth}}{
        \cellcolor[rgb]{0.9,0.9,0.9}{
            \makecell[{{p{\linewidth}}}]{
                \texttt{\tiny{[GM$\langle$P2]}}
                \texttt{BOG: D, (6,4)} \\
            }
        }
    }
     \\ \\

    \theutterance \stepcounter{utterance}  
    & & & \multicolumn{2}{p{0.3\linewidth}}{
        \cellcolor[rgb]{0.9,0.9,0.9}{
            \makecell[{{p{\linewidth}}}]{
                \texttt{\tiny{[GM$|$GM]}}
                \texttt{BOG: D, (6,4)} \\
            }
        }
    }
    & & \\ \\

    \theutterance \stepcounter{utterance}  
    & & & \multicolumn{2}{p{0.3\linewidth}}{
        \cellcolor[rgb]{0.9,0.9,0.9}{
            \makecell[{{p{\linewidth}}}]{
                \texttt{\tiny{[GM$|$GM]}}
                \texttt{Pionós: níl (6,4) folamh, ach tá '┌' ann.} \\
\\ 
\texttt{Tá 3 as 16 pionós carntha agaibh le chéile.} \\
\texttt{Déan iarracht arís, le do thoil!} \\
            }
        }
    }
    & & \\ \\

    \theutterance \stepcounter{utterance}  
    & & \multicolumn{4}{p{0.6\linewidth}}{
        \cellcolor[rgb]{0.9,0.9,0.9}{
            \makecell[{{p{\linewidth}}}]{
                \texttt{\tiny{[GM$\rangle$P2]}}
                \texttt{Pionós: níl (6,4) folamh, ach tá '┌' ann.} \\
\\ 
\texttt{Tá 3 as 16 pionós carntha agaibh le chéile.} \\
\texttt{Déan iarracht arís, le do thoil!} \\
            }
        }
    }
    & \\ \\

    \theutterance \stepcounter{utterance}  
    & & & \multicolumn{4}{p{0.6\linewidth}}{
        \cellcolor[rgb]{0.9,0.9,0.9}{
            \makecell[{{p{\linewidth}}}]{
                \texttt{\tiny{[GM$\langle$P2]}}
                \texttt{BOG: D, (5,5)} \\
            }
        }
    }
     \\ \\

    \theutterance \stepcounter{utterance}  
    & & & \multicolumn{2}{p{0.3\linewidth}}{
        \cellcolor[rgb]{0.9,0.9,0.9}{
            \makecell[{{p{\linewidth}}}]{
                \texttt{\tiny{[GM$|$GM]}}
                \texttt{BOG: D, (5,5)} \\
            }
        }
    }
    & & \\ \\

    \theutterance \stepcounter{utterance}  
    & & & \multicolumn{2}{p{0.3\linewidth}}{
        \cellcolor[rgb]{0.9,0.9,0.9}{
            \makecell[{{p{\linewidth}}}]{
                \texttt{\tiny{[GM$|$GM]}}
                \texttt{D moved to (5, 5)} \\
            }
        }
    }
    & & \\ \\

    \theutterance \stepcounter{utterance}  
    & & \multicolumn{4}{p{0.6\linewidth}}{
        \cellcolor[rgb]{0.9,0.9,0.9}{
            \makecell[{{p{\linewidth}}}]{
                \texttt{\tiny{[P1$\langle$GM]}}
                \texttt{Bogadh 'O' go (5,3) go rathúil. Seo mar atá do ghreille nuashonraithe:} \\
\texttt{```} \\
\texttt{ 1234567} \\
\texttt{╔═╤══╤══╗} \\
\texttt{║◌│◌T│◌◌║ 1} \\
\texttt{║◌└┬─┴──╢ 2} \\
\texttt{║◌◌│◌O◌◌║ 3} \\
\texttt{╟─┬┘◌P┌─╢ 4} \\
\texttt{║◌│AD◌│◌║ 5} \\
\texttt{╟─┼─┬─┤M║ 6} \\
\texttt{║◌│S│◌│◌║ 7} \\
\texttt{╚═╧═╧═╧═╝} \\
\texttt{```} \\
\texttt{Tá sibh ag imirt babhta 5 faoi láthair as uasmhéid 28.} \\
\texttt{Tá 3 as 16 pionós carntha agaibh le chéile.} \\
\texttt{Bhog an t{-}imreoir eile réad ar a ngreille.} \\
\\ 
\\ 
\texttt{Cad é do chéad ordú eile?} \\
            }
        }
    }
    & \\ \\

    \theutterance \stepcounter{utterance}  
    & \multicolumn{4}{p{0.6\linewidth}}{
        \cellcolor[rgb]{0.9,0.9,0.9}{
            \makecell[{{p{\linewidth}}}]{
                \texttt{\tiny{[P1$\rangle$GM]}}
                \texttt{BOG: T, (5,2)} \\
            }
        }
    }
    & & \\ \\

    \theutterance \stepcounter{utterance}  
    & & & \multicolumn{2}{p{0.3\linewidth}}{
        \cellcolor[rgb]{0.9,0.9,0.9}{
            \makecell[{{p{\linewidth}}}]{
                \texttt{\tiny{[GM$|$GM]}}
                \texttt{BOG: T, (5,2)} \\
            }
        }
    }
    & & \\ \\

    \theutterance \stepcounter{utterance}  
    & & & \multicolumn{2}{p{0.3\linewidth}}{
        \cellcolor[rgb]{0.9,0.9,0.9}{
            \makecell[{{p{\linewidth}}}]{
                \texttt{\tiny{[GM$|$GM]}}
                \texttt{Pionós: níl (5,2) folamh, ach tá '┴' ann.} \\
\\ 
\texttt{Tá 4 as 16 pionós carntha agaibh le chéile.} \\
\texttt{Déan iarracht arís, le do thoil!} \\
            }
        }
    }
    & & \\ \\

    \theutterance \stepcounter{utterance}  
    & & \multicolumn{4}{p{0.6\linewidth}}{
        \cellcolor[rgb]{0.9,0.9,0.9}{
            \makecell[{{p{\linewidth}}}]{
                \texttt{\tiny{[P1$\langle$GM]}}
                \texttt{Pionós: níl (5,2) folamh, ach tá '┴' ann.} \\
\\ 
\texttt{Tá 4 as 16 pionós carntha agaibh le chéile.} \\
\texttt{Déan iarracht arís, le do thoil!} \\
            }
        }
    }
    & \\ \\

    \theutterance \stepcounter{utterance}  
    & \multicolumn{4}{p{0.6\linewidth}}{
        \cellcolor[rgb]{0.9,0.9,0.9}{
            \makecell[{{p{\linewidth}}}]{
                \texttt{\tiny{[P1$\rangle$GM]}}
                \texttt{BOG: T, (6,2)} \\
            }
        }
    }
    & & \\ \\

    \theutterance \stepcounter{utterance}  
    & & & \multicolumn{2}{p{0.3\linewidth}}{
        \cellcolor[rgb]{0.9,0.9,0.9}{
            \makecell[{{p{\linewidth}}}]{
                \texttt{\tiny{[GM$|$GM]}}
                \texttt{BOG: T, (6,2)} \\
            }
        }
    }
    & & \\ \\

    \theutterance \stepcounter{utterance}  
    & & & \multicolumn{2}{p{0.3\linewidth}}{
        \cellcolor[rgb]{0.9,0.9,0.9}{
            \makecell[{{p{\linewidth}}}]{
                \texttt{\tiny{[GM$|$GM]}}
                \texttt{Pionós: níl (6,2) folamh, ach tá '─' ann.} \\
\\ 
\texttt{Tá 5 as 16 pionós carntha agaibh le chéile.} \\
\texttt{Déan iarracht arís, le do thoil!} \\
            }
        }
    }
    & & \\ \\

    \theutterance \stepcounter{utterance}  
    & & \multicolumn{4}{p{0.6\linewidth}}{
        \cellcolor[rgb]{0.9,0.9,0.9}{
            \makecell[{{p{\linewidth}}}]{
                \texttt{\tiny{[P1$\langle$GM]}}
                \texttt{Pionós: níl (6,2) folamh, ach tá '─' ann.} \\
\\ 
\texttt{Tá 5 as 16 pionós carntha agaibh le chéile.} \\
\texttt{Déan iarracht arís, le do thoil!} \\
            }
        }
    }
    & \\ \\

    \theutterance \stepcounter{utterance}  
    & \multicolumn{4}{p{0.6\linewidth}}{
        \cellcolor[rgb]{0.9,0.9,0.9}{
            \makecell[{{p{\linewidth}}}]{
                \texttt{\tiny{[P1$\rangle$GM]}}
                \texttt{BOG: T, (1,2)} \\
            }
        }
    }
    & & \\ \\

    \theutterance \stepcounter{utterance}  
    & & & \multicolumn{2}{p{0.3\linewidth}}{
        \cellcolor[rgb]{0.9,0.9,0.9}{
            \makecell[{{p{\linewidth}}}]{
                \texttt{\tiny{[GM$|$GM]}}
                \texttt{BOG: T, (1,2)} \\
            }
        }
    }
    & & \\ \\

    \theutterance \stepcounter{utterance}  
    & & & \multicolumn{2}{p{0.3\linewidth}}{
        \cellcolor[rgb]{0.9,0.9,0.9}{
            \makecell[{{p{\linewidth}}}]{
                \texttt{\tiny{[GM$|$GM]}}
                \texttt{T moved to (1, 2)} \\
            }
        }
    }
    & & \\ \\

    \theutterance \stepcounter{utterance}  
    & & \multicolumn{4}{p{0.6\linewidth}}{
        \cellcolor[rgb]{0.9,0.9,0.9}{
            \makecell[{{p{\linewidth}}}]{
                \texttt{\tiny{[GM$\rangle$P2]}}
                \texttt{Bogadh 'D' go (5,5) go rathúil. Seo mar atá do ghreille nuashonraithe:} \\
\texttt{```} \\
\texttt{ 1234567} \\
\texttt{╔═╤══╤══╗} \\
\texttt{║◌│O◌│◌◌║ 1} \\
\texttt{║◌└┬─┴──╢ 2} \\
\texttt{║◌◌│◌◌◌◌║ 3} \\
\texttt{╟─┬┘◌P┌─╢ 4} \\
\texttt{║◌│◌◌D│A║ 5} \\
\texttt{╟─┼─┬─┤S║ 6} \\
\texttt{║M│◌│◌│T║ 7} \\
\texttt{╚═╧═╧═╧═╝} \\
\texttt{```} \\
\texttt{Tá sibh ag imirt babhta 6 faoi láthair as uasmhéid 28.} \\
\texttt{Tá 5 as 16 pionós carntha agaibh le chéile.} \\
\texttt{Bhog an t{-}imreoir eile réad ar a ngreille.} \\
\\ 
\\ 
\texttt{Cad é do chéad ordú eile?} \\
            }
        }
    }
    & \\ \\

    \theutterance \stepcounter{utterance}  
    & & & \multicolumn{4}{p{0.6\linewidth}}{
        \cellcolor[rgb]{0.9,0.9,0.9}{
            \makecell[{{p{\linewidth}}}]{
                \texttt{\tiny{[GM$\langle$P2]}}
                \texttt{BOG: O, (5,3)} \\
            }
        }
    }
     \\ \\

    \theutterance \stepcounter{utterance}  
    & & & \multicolumn{2}{p{0.3\linewidth}}{
        \cellcolor[rgb]{0.9,0.9,0.9}{
            \makecell[{{p{\linewidth}}}]{
                \texttt{\tiny{[GM$|$GM]}}
                \texttt{BOG: O, (5,3)} \\
            }
        }
    }
    & & \\ \\

    \theutterance \stepcounter{utterance}  
    & & & \multicolumn{2}{p{0.3\linewidth}}{
        \cellcolor[rgb]{0.9,0.9,0.9}{
            \makecell[{{p{\linewidth}}}]{
                \texttt{\tiny{[GM$|$GM]}}
                \texttt{O moved to (5, 3)} \\
            }
        }
    }
    & & \\ \\

    \theutterance \stepcounter{utterance}  
    & & \multicolumn{4}{p{0.6\linewidth}}{
        \cellcolor[rgb]{0.9,0.9,0.9}{
            \makecell[{{p{\linewidth}}}]{
                \texttt{\tiny{[P1$\langle$GM]}}
                \texttt{Bogadh 'T' go (1,2) go rathúil. Seo mar atá do ghreille nuashonraithe:} \\
\texttt{```} \\
\texttt{ 1234567} \\
\texttt{╔═╤══╤══╗} \\
\texttt{║◌│◌◌│◌◌║ 1} \\
\texttt{║T└┬─┴──╢ 2} \\
\texttt{║◌◌│◌O◌◌║ 3} \\
\texttt{╟─┬┘◌P┌─╢ 4} \\
\texttt{║◌│AD◌│◌║ 5} \\
\texttt{╟─┼─┬─┤M║ 6} \\
\texttt{║◌│S│◌│◌║ 7} \\
\texttt{╚═╧═╧═╧═╝} \\
\texttt{```} \\
\texttt{Tá sibh ag imirt babhta 6 faoi láthair as uasmhéid 28.} \\
\texttt{Tá 5 as 16 pionós carntha agaibh le chéile.} \\
\texttt{Bhog an t{-}imreoir eile réad ar a ngreille.} \\
\\ 
\\ 
\texttt{Cad é do chéad ordú eile?} \\
            }
        }
    }
    & \\ \\

    \theutterance \stepcounter{utterance}  
    & \multicolumn{4}{p{0.6\linewidth}}{
        \cellcolor[rgb]{0.9,0.9,0.9}{
            \makecell[{{p{\linewidth}}}]{
                \texttt{\tiny{[P1$\rangle$GM]}}
                \texttt{ABAIR: Bhog mé P go (5,4) agus O go (5,3). Ní raibh mé in ann T a bhogadh go (5,2) nó (6,2) mar gheall ar línte, mar sin bhog mé T go (1,2) go sealadach. An bhfuil sprioc eile againn do T? B'fhéidir (7,2) nó (6,3)? Nó an féidir leat T a chur in áit éigin eile agus mé ag cur mo cheannsa ann? Cad iad na sprioc{-}shuímh is fearr duitse anois?} \\
            }
        }
    }
    & & \\ \\

    \theutterance \stepcounter{utterance}  
    & & & \multicolumn{2}{p{0.3\linewidth}}{
        \cellcolor[rgb]{0.9,0.9,0.9}{
            \makecell[{{p{\linewidth}}}]{
                \texttt{\tiny{[GM$|$GM]}}
                \texttt{ABAIR: Bhog mé P go (5,4) agus O go (5,3). Ní raibh mé in ann T a bhogadh go (5,2) nó (6,2) mar gheall ar línte, mar sin bhog mé T go (1,2) go sealadach. An bhfuil sprioc eile againn do T? B'fhéidir (7,2) nó (6,3)? Nó an féidir leat T a chur in áit éigin eile agus mé ag cur mo cheannsa ann? Cad iad na sprioc{-}shuímh is fearr duitse anois?} \\
            }
        }
    }
    & & \\ \\

    \theutterance \stepcounter{utterance}  
    & & \multicolumn{4}{p{0.6\linewidth}}{
        \cellcolor[rgb]{0.9,0.9,0.9}{
            \makecell[{{p{\linewidth}}}]{
                \texttt{\tiny{[GM$\rangle$P2]}}
                \texttt{Bogadh 'O' go (5,3) go rathúil. Seo mar atá do ghreille nuashonraithe:} \\
\texttt{```} \\
\texttt{ 1234567} \\
\texttt{╔═╤══╤══╗} \\
\texttt{║◌│◌◌│◌◌║ 1} \\
\texttt{║◌└┬─┴──╢ 2} \\
\texttt{║◌◌│◌O◌◌║ 3} \\
\texttt{╟─┬┘◌P┌─╢ 4} \\
\texttt{║◌│◌◌D│A║ 5} \\
\texttt{╟─┼─┬─┤S║ 6} \\
\texttt{║M│◌│◌│T║ 7} \\
\texttt{╚═╧═╧═╧═╝} \\
\texttt{```} \\
\texttt{Tá sibh ag imirt babhta 7 faoi láthair as uasmhéid 28.} \\
\texttt{Tá 5 as 16 pionós carntha agaibh le chéile.} \\
\texttt{Sheol an t{-}imreoir eile an teachtaireacht seo:} \\
\\ 
\\ 
\texttt{"Bhog mé P go (5,4) agus O go (5,3). Ní raibh mé in ann T a bhogadh go (5,2) nó (6,2) mar gheall ar línte, mar sin bhog mé T go (1,2) go sealadach. An bhfuil sprioc eile againn do T? B'fhéidir (7,2) nó (6,3)? Nó an féidir leat T a chur in áit éigin eile agus mé ag cur mo cheannsa ann? Cad iad na sprioc{-}shuímh is fearr duitse anois?"} \\
\\ 
\\ 
\texttt{Cad é do chéad ordú eile?} \\
            }
        }
    }
    & \\ \\

    \theutterance \stepcounter{utterance}  
    & & & \multicolumn{4}{p{0.6\linewidth}}{
        \cellcolor[rgb]{0.9,0.9,0.9}{
            \makecell[{{p{\linewidth}}}]{
                \texttt{\tiny{[GM$\langle$P2]}}
                \texttt{ABAIR: Go raibh maith agat as an nuashonrú. Seiceáil mé mo chlár: (5,2) folamh, (6,3) folamh, (7,2) tá fráma ann. Is féidir liom T a bhogadh go (5,2) nó (6,3). Ar mhaith leat T a chur ag (5,2) mar a bhí sa phlean bunaidh, nó ag (6,3)? Faoi láthair tá mo T ag (7,7) agus is féidir liom é a bhogadh. D'éirigh liom D a bhogadh go (5,5) agus O go (5,3). Bogfaidh mé S agus A go luath. Cad a cheapann tú?} \\
            }
        }
    }
     \\ \\

    \theutterance \stepcounter{utterance}  
    & & & \multicolumn{2}{p{0.3\linewidth}}{
        \cellcolor[rgb]{0.9,0.9,0.9}{
            \makecell[{{p{\linewidth}}}]{
                \texttt{\tiny{[GM$|$GM]}}
                \texttt{ABAIR: Go raibh maith agat as an nuashonrú. Seiceáil mé mo chlár: (5,2) folamh, (6,3) folamh, (7,2) tá fráma ann. Is féidir liom T a bhogadh go (5,2) nó (6,3). Ar mhaith leat T a chur ag (5,2) mar a bhí sa phlean bunaidh, nó ag (6,3)? Faoi láthair tá mo T ag (7,7) agus is féidir liom é a bhogadh. D'éirigh liom D a bhogadh go (5,5) agus O go (5,3). Bogfaidh mé S agus A go luath. Cad a cheapann tú?} \\
            }
        }
    }
    & & \\ \\

    \theutterance \stepcounter{utterance}  
    & & \multicolumn{4}{p{0.6\linewidth}}{
        \cellcolor[rgb]{0.9,0.9,0.9}{
            \makecell[{{p{\linewidth}}}]{
                \texttt{\tiny{[P1$\langle$GM]}}
                \texttt{Cuireadh do theachtaireacht ar aghaidh chuig an imreoir eile.} \\
\\ 
\\ 
\texttt{Tá sibh ag imirt babhta 7 faoi láthair as uasmhéid 28.} \\
\texttt{Tá 5 as 16 pionós carntha agaibh le chéile.} \\
\texttt{Sheol an t{-}imreoir eile an teachtaireacht seo:} \\
\\ 
\\ 
\texttt{"Go raibh maith agat as an nuashonrú. Seiceáil mé mo chlár: (5,2) folamh, (6,3) folamh, (7,2) tá fráma ann. Is féidir liom T a bhogadh go (5,2) nó (6,3). Ar mhaith leat T a chur ag (5,2) mar a bhí sa phlean bunaidh, nó ag (6,3)? Faoi láthair tá mo T ag (7,7) agus is féidir liom é a bhogadh. D'éirigh liom D a bhogadh go (5,5) agus O go (5,3). Bogfaidh mé S agus A go luath. Cad a cheapann tú?"} \\
\\ 
\\ 
\texttt{Cad é do chéad ordú eile?} \\
            }
        }
    }
    & \\ \\

    \theutterance \stepcounter{utterance}  
    & \multicolumn{4}{p{0.6\linewidth}}{
        \cellcolor[rgb]{0.9,0.9,0.9}{
            \makecell[{{p{\linewidth}}}]{
                \texttt{\tiny{[P1$\rangle$GM]}}
                \texttt{ABAIR: Is é (5,2) an sprioc bunaidh ach tá líne ann ar mo chlár, mar sin ní féidir liom T a chur ann. Bain triail as (6,3) do T. Bogfaidh mé mo T ó (1,2) go (6,3) nuair a bheidh sé folamh. Idir an dá linn, bogfaidh mé S go (4,4) agus A go (5,5) de réir an phlean. An bhfuil (5,5) folamh agat? Tá D agam cheana ag (3,5) agus caithfidh mé é a bhogadh go (6,4).} \\
            }
        }
    }
    & & \\ \\

    \theutterance \stepcounter{utterance}  
    & & & \multicolumn{2}{p{0.3\linewidth}}{
        \cellcolor[rgb]{0.9,0.9,0.9}{
            \makecell[{{p{\linewidth}}}]{
                \texttt{\tiny{[GM$|$GM]}}
                \texttt{ABAIR: Is é (5,2) an sprioc bunaidh ach tá líne ann ar mo chlár, mar sin ní féidir liom T a chur ann. Bain triail as (6,3) do T. Bogfaidh mé mo T ó (1,2) go (6,3) nuair a bheidh sé folamh. Idir an dá linn, bogfaidh mé S go (4,4) agus A go (5,5) de réir an phlean. An bhfuil (5,5) folamh agat? Tá D agam cheana ag (3,5) agus caithfidh mé é a bhogadh go (6,4).} \\
            }
        }
    }
    & & \\ \\

    \theutterance \stepcounter{utterance}  
    & & \multicolumn{4}{p{0.6\linewidth}}{
        \cellcolor[rgb]{0.9,0.9,0.9}{
            \makecell[{{p{\linewidth}}}]{
                \texttt{\tiny{[GM$\rangle$P2]}}
                \texttt{Cuireadh do theachtaireacht ar aghaidh chuig an imreoir eile.} \\
\\ 
\\ 
\texttt{Tá sibh ag imirt babhta 8 faoi láthair as uasmhéid 28.} \\
\texttt{Tá 5 as 16 pionós carntha agaibh le chéile.} \\
\texttt{Sheol an t{-}imreoir eile an teachtaireacht seo:} \\
\\ 
\\ 
\texttt{"Is é (5,2) an sprioc bunaidh ach tá líne ann ar mo chlár, mar sin ní féidir liom T a chur ann. Bain triail as (6,3) do T. Bogfaidh mé mo T ó (1,2) go (6,3) nuair a bheidh sé folamh. Idir an dá linn, bogfaidh mé S go (4,4) agus A go (5,5) de réir an phlean. An bhfuil (5,5) folamh agat? Tá D agam cheana ag (3,5) agus caithfidh mé é a bhogadh go (6,4)."} \\
\\ 
\\ 
\texttt{Cad é do chéad ordú eile?} \\
            }
        }
    }
    & \\ \\

    \theutterance \stepcounter{utterance}  
    & & & \multicolumn{4}{p{0.6\linewidth}}{
        \cellcolor[rgb]{0.9,0.9,0.9}{
            \makecell[{{p{\linewidth}}}]{
                \texttt{\tiny{[GM$\langle$P2]}}
                \texttt{ABAIR: Tá D agam cheana ag (5,5), mar sin níl (5,5) folamh. Ní mór dúinn A a chur in áit éigin eile. An féidir linn A a chur ag (5,6) nó (6,5) ina ionad? Seiceáil mé mo chlár: (5,6) folamh, (6,5) folamh. Bogfaidh mé T go (6,3) anois. Maidir le S ag (4,4), tá sé folamh ar mo chlár, mar sin is féidir liom S a bhogadh ann. Ach beidh orainn A a choigeartú. Cad é do rogha?} \\
            }
        }
    }
     \\ \\

    \theutterance \stepcounter{utterance}  
    & & & \multicolumn{2}{p{0.3\linewidth}}{
        \cellcolor[rgb]{0.9,0.9,0.9}{
            \makecell[{{p{\linewidth}}}]{
                \texttt{\tiny{[GM$|$GM]}}
                \texttt{ABAIR: Tá D agam cheana ag (5,5), mar sin níl (5,5) folamh. Ní mór dúinn A a chur in áit éigin eile. An féidir linn A a chur ag (5,6) nó (6,5) ina ionad? Seiceáil mé mo chlár: (5,6) folamh, (6,5) folamh. Bogfaidh mé T go (6,3) anois. Maidir le S ag (4,4), tá sé folamh ar mo chlár, mar sin is féidir liom S a bhogadh ann. Ach beidh orainn A a choigeartú. Cad é do rogha?} \\
            }
        }
    }
    & & \\ \\

    \theutterance \stepcounter{utterance}  
    & & \multicolumn{4}{p{0.6\linewidth}}{
        \cellcolor[rgb]{0.9,0.9,0.9}{
            \makecell[{{p{\linewidth}}}]{
                \texttt{\tiny{[P1$\langle$GM]}}
                \texttt{Cuireadh do theachtaireacht ar aghaidh chuig an imreoir eile.} \\
\\ 
\\ 
\texttt{Tá sibh ag imirt babhta 8 faoi láthair as uasmhéid 28.} \\
\texttt{Tá 5 as 16 pionós carntha agaibh le chéile.} \\
\texttt{Sheol an t{-}imreoir eile an teachtaireacht seo:} \\
\\ 
\\ 
\texttt{"Tá D agam cheana ag (5,5), mar sin níl (5,5) folamh. Ní mór dúinn A a chur in áit éigin eile. An féidir linn A a chur ag (5,6) nó (6,5) ina ionad? Seiceáil mé mo chlár: (5,6) folamh, (6,5) folamh. Bogfaidh mé T go (6,3) anois. Maidir le S ag (4,4), tá sé folamh ar mo chlár, mar sin is féidir liom S a bhogadh ann. Ach beidh orainn A a choigeartú. Cad é do rogha?"} \\
\\ 
\\ 
\texttt{Cad é do chéad ordú eile?} \\
            }
        }
    }
    & \\ \\

    \theutterance \stepcounter{utterance}  
    & \multicolumn{4}{p{0.6\linewidth}}{
        \cellcolor[rgb]{0.9,0.9,0.9}{
            \makecell[{{p{\linewidth}}}]{
                \texttt{\tiny{[P1$\rangle$GM]}}
                \texttt{ABAIR: Roghnaigh muid A ag (5,6) mar sin. Tá (5,6) folamh ar mo chlár. Bogfaidh mé S go (4,4) anois. Tar éis sin, bogfaidh mé D ó (3,5) go (6,4), agus A ó (2,5) go (5,6). Ansin nuair a bheidh T bogtha agat go (6,3), bogfaidh mé mo T ó (1,2) go (6,3). Leanfaimid orainn!} \\
            }
        }
    }
    & & \\ \\

    \theutterance \stepcounter{utterance}  
    & & & \multicolumn{2}{p{0.3\linewidth}}{
        \cellcolor[rgb]{0.9,0.9,0.9}{
            \makecell[{{p{\linewidth}}}]{
                \texttt{\tiny{[GM$|$GM]}}
                \texttt{ABAIR: Roghnaigh muid A ag (5,6) mar sin. Tá (5,6) folamh ar mo chlár. Bogfaidh mé S go (4,4) anois. Tar éis sin, bogfaidh mé D ó (3,5) go (6,4), agus A ó (2,5) go (5,6). Ansin nuair a bheidh T bogtha agat go (6,3), bogfaidh mé mo T ó (1,2) go (6,3). Leanfaimid orainn!} \\
            }
        }
    }
    & & \\ \\

    \theutterance \stepcounter{utterance}  
    & & \multicolumn{4}{p{0.6\linewidth}}{
        \cellcolor[rgb]{0.9,0.9,0.9}{
            \makecell[{{p{\linewidth}}}]{
                \texttt{\tiny{[GM$\rangle$P2]}}
                \texttt{Cuireadh do theachtaireacht ar aghaidh chuig an imreoir eile.} \\
\\ 
\\ 
\texttt{Tá sibh ag imirt babhta 9 faoi láthair as uasmhéid 28.} \\
\texttt{Tá 5 as 16 pionós carntha agaibh le chéile.} \\
\texttt{Sheol an t{-}imreoir eile an teachtaireacht seo:} \\
\\ 
\\ 
\texttt{"Roghnaigh muid A ag (5,6) mar sin. Tá (5,6) folamh ar mo chlár. Bogfaidh mé S go (4,4) anois. Tar éis sin, bogfaidh mé D ó (3,5) go (6,4), agus A ó (2,5) go (5,6). Ansin nuair a bheidh T bogtha agat go (6,3), bogfaidh mé mo T ó (1,2) go (6,3). Leanfaimid orainn!"} \\
\\ 
\\ 
\texttt{Cad é do chéad ordú eile?} \\
            }
        }
    }
    & \\ \\

    \theutterance \stepcounter{utterance}  
    & & & \multicolumn{4}{p{0.6\linewidth}}{
        \cellcolor[rgb]{0.9,0.9,0.9}{
            \makecell[{{p{\linewidth}}}]{
                \texttt{\tiny{[GM$\langle$P2]}}
                \texttt{BOG: T, (6,3)} \\
            }
        }
    }
     \\ \\

    \theutterance \stepcounter{utterance}  
    & & & \multicolumn{2}{p{0.3\linewidth}}{
        \cellcolor[rgb]{0.9,0.9,0.9}{
            \makecell[{{p{\linewidth}}}]{
                \texttt{\tiny{[GM$|$GM]}}
                \texttt{BOG: T, (6,3)} \\
            }
        }
    }
    & & \\ \\

    \theutterance \stepcounter{utterance}  
    & & & \multicolumn{2}{p{0.3\linewidth}}{
        \cellcolor[rgb]{0.9,0.9,0.9}{
            \makecell[{{p{\linewidth}}}]{
                \texttt{\tiny{[GM$|$GM]}}
                \texttt{T moved to (6, 3)} \\
            }
        }
    }
    & & \\ \\

    \theutterance \stepcounter{utterance}  
    & & \multicolumn{4}{p{0.6\linewidth}}{
        \cellcolor[rgb]{0.9,0.9,0.9}{
            \makecell[{{p{\linewidth}}}]{
                \texttt{\tiny{[P1$\langle$GM]}}
                \texttt{Cuireadh do theachtaireacht ar aghaidh chuig an imreoir eile.} \\
\\ 
\\ 
\texttt{Tá sibh ag imirt babhta 9 faoi láthair as uasmhéid 28.} \\
\texttt{Tá 5 as 16 pionós carntha agaibh le chéile.} \\
\texttt{Bhog an t{-}imreoir eile réad ar a ngreille.} \\
\\ 
\\ 
\texttt{Cad é do chéad ordú eile?} \\
            }
        }
    }
    & \\ \\

    \theutterance \stepcounter{utterance}  
    & \multicolumn{4}{p{0.6\linewidth}}{
        \cellcolor[rgb]{0.9,0.9,0.9}{
            \makecell[{{p{\linewidth}}}]{
                \texttt{\tiny{[P1$\rangle$GM]}}
                \texttt{BOG: S, (4,4)} \\
            }
        }
    }
    & & \\ \\

    \theutterance \stepcounter{utterance}  
    & & & \multicolumn{2}{p{0.3\linewidth}}{
        \cellcolor[rgb]{0.9,0.9,0.9}{
            \makecell[{{p{\linewidth}}}]{
                \texttt{\tiny{[GM$|$GM]}}
                \texttt{BOG: S, (4,4)} \\
            }
        }
    }
    & & \\ \\

    \theutterance \stepcounter{utterance}  
    & & & \multicolumn{2}{p{0.3\linewidth}}{
        \cellcolor[rgb]{0.9,0.9,0.9}{
            \makecell[{{p{\linewidth}}}]{
                \texttt{\tiny{[GM$|$GM]}}
                \texttt{S moved to (4, 4)} \\
            }
        }
    }
    & & \\ \\

    \theutterance \stepcounter{utterance}  
    & & \multicolumn{4}{p{0.6\linewidth}}{
        \cellcolor[rgb]{0.9,0.9,0.9}{
            \makecell[{{p{\linewidth}}}]{
                \texttt{\tiny{[GM$\rangle$P2]}}
                \texttt{Bogadh 'T' go (6,3) go rathúil. Seo mar atá do ghreille nuashonraithe:} \\
\texttt{```} \\
\texttt{ 1234567} \\
\texttt{╔═╤══╤══╗} \\
\texttt{║◌│◌◌│◌◌║ 1} \\
\texttt{║◌└┬─┴──╢ 2} \\
\texttt{║◌◌│◌OT◌║ 3} \\
\texttt{╟─┬┘◌P┌─╢ 4} \\
\texttt{║◌│◌◌D│A║ 5} \\
\texttt{╟─┼─┬─┤S║ 6} \\
\texttt{║M│◌│◌│◌║ 7} \\
\texttt{╚═╧═╧═╧═╝} \\
\texttt{```} \\
\texttt{Tá sibh ag imirt babhta 10 faoi láthair as uasmhéid 28.} \\
\texttt{Tá 5 as 16 pionós carntha agaibh le chéile.} \\
\texttt{Bhog an t{-}imreoir eile réad ar a ngreille.} \\
\\ 
\\ 
\texttt{Cad é do chéad ordú eile?} \\
            }
        }
    }
    & \\ \\

    \theutterance \stepcounter{utterance}  
    & & & \multicolumn{4}{p{0.6\linewidth}}{
        \cellcolor[rgb]{0.9,0.9,0.9}{
            \makecell[{{p{\linewidth}}}]{
                \texttt{\tiny{[GM$\langle$P2]}}
                \texttt{BOG: S, (4,4)} \\
            }
        }
    }
     \\ \\

    \theutterance \stepcounter{utterance}  
    & & & \multicolumn{2}{p{0.3\linewidth}}{
        \cellcolor[rgb]{0.9,0.9,0.9}{
            \makecell[{{p{\linewidth}}}]{
                \texttt{\tiny{[GM$|$GM]}}
                \texttt{BOG: S, (4,4)} \\
            }
        }
    }
    & & \\ \\

    \theutterance \stepcounter{utterance}  
    & & & \multicolumn{2}{p{0.3\linewidth}}{
        \cellcolor[rgb]{0.9,0.9,0.9}{
            \makecell[{{p{\linewidth}}}]{
                \texttt{\tiny{[GM$|$GM]}}
                \texttt{S moved to (4, 4)} \\
            }
        }
    }
    & & \\ \\

    \theutterance \stepcounter{utterance}  
    & & \multicolumn{4}{p{0.6\linewidth}}{
        \cellcolor[rgb]{0.9,0.9,0.9}{
            \makecell[{{p{\linewidth}}}]{
                \texttt{\tiny{[P1$\langle$GM]}}
                \texttt{Bogadh 'S' go (4,4) go rathúil. Seo mar atá do ghreille nuashonraithe:} \\
\texttt{```} \\
\texttt{ 1234567} \\
\texttt{╔═╤══╤══╗} \\
\texttt{║◌│◌◌│◌◌║ 1} \\
\texttt{║T└┬─┴──╢ 2} \\
\texttt{║◌◌│◌O◌◌║ 3} \\
\texttt{╟─┬┘SP┌─╢ 4} \\
\texttt{║◌│AD◌│◌║ 5} \\
\texttt{╟─┼─┬─┤M║ 6} \\
\texttt{║◌│◌│◌│◌║ 7} \\
\texttt{╚═╧═╧═╧═╝} \\
\texttt{```} \\
\texttt{Tá sibh ag imirt babhta 10 faoi láthair as uasmhéid 28.} \\
\texttt{Tá 5 as 16 pionós carntha agaibh le chéile.} \\
\texttt{Bhog an t{-}imreoir eile réad ar a ngreille.} \\
\\ 
\\ 
\texttt{Cad é do chéad ordú eile?} \\
            }
        }
    }
    & \\ \\

    \theutterance \stepcounter{utterance}  
    & \multicolumn{4}{p{0.6\linewidth}}{
        \cellcolor[rgb]{0.9,0.9,0.9}{
            \makecell[{{p{\linewidth}}}]{
                \texttt{\tiny{[P1$\rangle$GM]}}
                \texttt{BOG: D, (6,4)} \\
            }
        }
    }
    & & \\ \\

    \theutterance \stepcounter{utterance}  
    & & & \multicolumn{2}{p{0.3\linewidth}}{
        \cellcolor[rgb]{0.9,0.9,0.9}{
            \makecell[{{p{\linewidth}}}]{
                \texttt{\tiny{[GM$|$GM]}}
                \texttt{BOG: D, (6,4)} \\
            }
        }
    }
    & & \\ \\

    \theutterance \stepcounter{utterance}  
    & & & \multicolumn{2}{p{0.3\linewidth}}{
        \cellcolor[rgb]{0.9,0.9,0.9}{
            \makecell[{{p{\linewidth}}}]{
                \texttt{\tiny{[GM$|$GM]}}
                \texttt{Pionós: níl (6,4) folamh, ach tá '┌' ann.} \\
\\ 
\texttt{Tá 6 as 16 pionós carntha agaibh le chéile.} \\
\texttt{Déan iarracht arís, le do thoil!} \\
            }
        }
    }
    & & \\ \\

    \theutterance \stepcounter{utterance}  
    & & \multicolumn{4}{p{0.6\linewidth}}{
        \cellcolor[rgb]{0.9,0.9,0.9}{
            \makecell[{{p{\linewidth}}}]{
                \texttt{\tiny{[P1$\langle$GM]}}
                \texttt{Pionós: níl (6,4) folamh, ach tá '┌' ann.} \\
\\ 
\texttt{Tá 6 as 16 pionós carntha agaibh le chéile.} \\
\texttt{Déan iarracht arís, le do thoil!} \\
            }
        }
    }
    & \\ \\

    \theutterance \stepcounter{utterance}  
    & \multicolumn{4}{p{0.6\linewidth}}{
        \cellcolor[rgb]{0.9,0.9,0.9}{
            \makecell[{{p{\linewidth}}}]{
                \texttt{\tiny{[P1$\rangle$GM]}}
                \texttt{BOG: D, (6,5)} \\
            }
        }
    }
    & & \\ \\

    \theutterance \stepcounter{utterance}  
    & & & \multicolumn{2}{p{0.3\linewidth}}{
        \cellcolor[rgb]{0.9,0.9,0.9}{
            \makecell[{{p{\linewidth}}}]{
                \texttt{\tiny{[GM$|$GM]}}
                \texttt{BOG: D, (6,5)} \\
            }
        }
    }
    & & \\ \\

    \theutterance \stepcounter{utterance}  
    & & & \multicolumn{2}{p{0.3\linewidth}}{
        \cellcolor[rgb]{0.9,0.9,0.9}{
            \makecell[{{p{\linewidth}}}]{
                \texttt{\tiny{[GM$|$GM]}}
                \texttt{Pionós: níl (6,5) folamh, ach tá '│' ann.} \\
\\ 
\texttt{Tá 7 as 16 pionós carntha agaibh le chéile.} \\
\texttt{Déan iarracht arís, le do thoil!} \\
            }
        }
    }
    & & \\ \\

    \theutterance \stepcounter{utterance}  
    & & \multicolumn{4}{p{0.6\linewidth}}{
        \cellcolor[rgb]{0.9,0.9,0.9}{
            \makecell[{{p{\linewidth}}}]{
                \texttt{\tiny{[P1$\langle$GM]}}
                \texttt{Pionós: níl (6,5) folamh, ach tá '│' ann.} \\
\\ 
\texttt{Tá 7 as 16 pionós carntha agaibh le chéile.} \\
\texttt{Déan iarracht arís, le do thoil!} \\
            }
        }
    }
    & \\ \\

    \theutterance \stepcounter{utterance}  
    & \multicolumn{4}{p{0.6\linewidth}}{
        \cellcolor[rgb]{0.9,0.9,0.9}{
            \makecell[{{p{\linewidth}}}]{
                \texttt{\tiny{[P1$\rangle$GM]}}
                \texttt{BOG: D, (7,5)} \\
            }
        }
    }
    & & \\ \\

    \theutterance \stepcounter{utterance}  
    & & & \multicolumn{2}{p{0.3\linewidth}}{
        \cellcolor[rgb]{0.9,0.9,0.9}{
            \makecell[{{p{\linewidth}}}]{
                \texttt{\tiny{[GM$|$GM]}}
                \texttt{BOG: D, (7,5)} \\
            }
        }
    }
    & & \\ \\

    \theutterance \stepcounter{utterance}  
    & & & \multicolumn{2}{p{0.3\linewidth}}{
        \cellcolor[rgb]{0.9,0.9,0.9}{
            \makecell[{{p{\linewidth}}}]{
                \texttt{\tiny{[GM$|$GM]}}
                \texttt{D moved to (7, 5)} \\
            }
        }
    }
    & & \\ \\

    \theutterance \stepcounter{utterance}  
    & & \multicolumn{4}{p{0.6\linewidth}}{
        \cellcolor[rgb]{0.9,0.9,0.9}{
            \makecell[{{p{\linewidth}}}]{
                \texttt{\tiny{[GM$\rangle$P2]}}
                \texttt{Bogadh 'S' go (4,4) go rathúil. Seo mar atá do ghreille nuashonraithe:} \\
\texttt{```} \\
\texttt{ 1234567} \\
\texttt{╔═╤══╤══╗} \\
\texttt{║◌│◌◌│◌◌║ 1} \\
\texttt{║◌└┬─┴──╢ 2} \\
\texttt{║◌◌│◌OT◌║ 3} \\
\texttt{╟─┬┘SP┌─╢ 4} \\
\texttt{║◌│◌◌D│A║ 5} \\
\texttt{╟─┼─┬─┤◌║ 6} \\
\texttt{║M│◌│◌│◌║ 7} \\
\texttt{╚═╧═╧═╧═╝} \\
\texttt{```} \\
\texttt{Tá sibh ag imirt babhta 11 faoi láthair as uasmhéid 28.} \\
\texttt{Tá 7 as 16 pionós carntha agaibh le chéile.} \\
\texttt{Bhog an t{-}imreoir eile réad ar a ngreille.} \\
\\ 
\\ 
\texttt{Cad é do chéad ordú eile?} \\
            }
        }
    }
    & \\ \\

    \theutterance \stepcounter{utterance}  
    & & & \multicolumn{4}{p{0.6\linewidth}}{
        \cellcolor[rgb]{0.9,0.9,0.9}{
            \makecell[{{p{\linewidth}}}]{
                \texttt{\tiny{[GM$\langle$P2]}}
                \texttt{BOG: A, (5,6)} \\
            }
        }
    }
     \\ \\

    \theutterance \stepcounter{utterance}  
    & & & \multicolumn{2}{p{0.3\linewidth}}{
        \cellcolor[rgb]{0.9,0.9,0.9}{
            \makecell[{{p{\linewidth}}}]{
                \texttt{\tiny{[GM$|$GM]}}
                \texttt{BOG: A, (5,6)} \\
            }
        }
    }
    & & \\ \\

    \theutterance \stepcounter{utterance}  
    & & & \multicolumn{2}{p{0.3\linewidth}}{
        \cellcolor[rgb]{0.9,0.9,0.9}{
            \makecell[{{p{\linewidth}}}]{
                \texttt{\tiny{[GM$|$GM]}}
                \texttt{Pionós: níl (5,6) folamh, ach tá '─' ann.} \\
\\ 
\texttt{Tá 8 as 16 pionós carntha agaibh le chéile.} \\
\texttt{Déan iarracht arís, le do thoil!} \\
            }
        }
    }
    & & \\ \\

    \theutterance \stepcounter{utterance}  
    & & \multicolumn{4}{p{0.6\linewidth}}{
        \cellcolor[rgb]{0.9,0.9,0.9}{
            \makecell[{{p{\linewidth}}}]{
                \texttt{\tiny{[GM$\rangle$P2]}}
                \texttt{Pionós: níl (5,6) folamh, ach tá '─' ann.} \\
\\ 
\texttt{Tá 8 as 16 pionós carntha agaibh le chéile.} \\
\texttt{Déan iarracht arís, le do thoil!} \\
            }
        }
    }
    & \\ \\

    \theutterance \stepcounter{utterance}  
    & & & \multicolumn{4}{p{0.6\linewidth}}{
        \cellcolor[rgb]{0.9,0.9,0.9}{
            \makecell[{{p{\linewidth}}}]{
                \texttt{\tiny{[GM$\langle$P2]}}
                \texttt{BOG: A, (6,5)} \\
            }
        }
    }
     \\ \\

    \theutterance \stepcounter{utterance}  
    & & & \multicolumn{2}{p{0.3\linewidth}}{
        \cellcolor[rgb]{0.9,0.9,0.9}{
            \makecell[{{p{\linewidth}}}]{
                \texttt{\tiny{[GM$|$GM]}}
                \texttt{BOG: A, (6,5)} \\
            }
        }
    }
    & & \\ \\

    \theutterance \stepcounter{utterance}  
    & & & \multicolumn{2}{p{0.3\linewidth}}{
        \cellcolor[rgb]{0.9,0.9,0.9}{
            \makecell[{{p{\linewidth}}}]{
                \texttt{\tiny{[GM$|$GM]}}
                \texttt{Pionós: níl (6,5) folamh, ach tá '│' ann.} \\
\\ 
\texttt{Tá 9 as 16 pionós carntha agaibh le chéile.} \\
\texttt{Déan iarracht arís, le do thoil!} \\
            }
        }
    }
    & & \\ \\

    \theutterance \stepcounter{utterance}  
    & & \multicolumn{4}{p{0.6\linewidth}}{
        \cellcolor[rgb]{0.9,0.9,0.9}{
            \makecell[{{p{\linewidth}}}]{
                \texttt{\tiny{[GM$\rangle$P2]}}
                \texttt{Pionós: níl (6,5) folamh, ach tá '│' ann.} \\
\\ 
\texttt{Tá 9 as 16 pionós carntha agaibh le chéile.} \\
\texttt{Déan iarracht arís, le do thoil!} \\
            }
        }
    }
    & \\ \\

    \theutterance \stepcounter{utterance}  
    & & & \multicolumn{4}{p{0.6\linewidth}}{
        \cellcolor[rgb]{0.9,0.9,0.9}{
            \makecell[{{p{\linewidth}}}]{
                \texttt{\tiny{[GM$\langle$P2]}}
                \texttt{BOG: A, (7,4)} \\
            }
        }
    }
     \\ \\

    \theutterance \stepcounter{utterance}  
    & & & \multicolumn{2}{p{0.3\linewidth}}{
        \cellcolor[rgb]{0.9,0.9,0.9}{
            \makecell[{{p{\linewidth}}}]{
                \texttt{\tiny{[GM$|$GM]}}
                \texttt{BOG: A, (7,4)} \\
            }
        }
    }
    & & \\ \\

    \theutterance \stepcounter{utterance}  
    & & & \multicolumn{2}{p{0.3\linewidth}}{
        \cellcolor[rgb]{0.9,0.9,0.9}{
            \makecell[{{p{\linewidth}}}]{
                \texttt{\tiny{[GM$|$GM]}}
                \texttt{Pionós: níl (7,4) folamh, ach tá '─' ann.} \\
\\ 
\texttt{Tá 10 as 16 pionós carntha agaibh le chéile.} \\
\texttt{Déan iarracht arís, le do thoil!} \\
            }
        }
    }
    & & \\ \\

    \theutterance \stepcounter{utterance}  
    & & \multicolumn{4}{p{0.6\linewidth}}{
        \cellcolor[rgb]{0.9,0.9,0.9}{
            \makecell[{{p{\linewidth}}}]{
                \texttt{\tiny{[GM$\rangle$P2]}}
                \texttt{Pionós: níl (7,4) folamh, ach tá '─' ann.} \\
\\ 
\texttt{Tá 10 as 16 pionós carntha agaibh le chéile.} \\
\texttt{Déan iarracht arís, le do thoil!} \\
            }
        }
    }
    & \\ \\

    \theutterance \stepcounter{utterance}  
    & & & \multicolumn{4}{p{0.6\linewidth}}{
        \cellcolor[rgb]{0.9,0.9,0.9}{
            \makecell[{{p{\linewidth}}}]{
                \texttt{\tiny{[GM$\langle$P2]}}
                \texttt{ABAIR: Tá deacracht agam A a bhogadh. Ar mo chlár, tá A ag (7,5) faoi láthair. Tá (5,6) blocáilte ag '─', (6,5) ag '│', agus (7,4) ag '─'. An féidir linn sprioc eile a roghnú do A? Seiceáil mé na cealla folmha in aice láimhe: (7,6) folamh, (6,6) folamh, (5,7) folamh, (4,6) folamh. An bhfuil ceann ar bith díobh seo oiriúnach duitse?} \\
            }
        }
    }
     \\ \\

    \theutterance \stepcounter{utterance}  
    & & & \multicolumn{2}{p{0.3\linewidth}}{
        \cellcolor[rgb]{0.9,0.9,0.9}{
            \makecell[{{p{\linewidth}}}]{
                \texttt{\tiny{[GM$|$GM]}}
                \texttt{ABAIR: Tá deacracht agam A a bhogadh. Ar mo chlár, tá A ag (7,5) faoi láthair. Tá (5,6) blocáilte ag '─', (6,5) ag '│', agus (7,4) ag '─'. An féidir linn sprioc eile a roghnú do A? Seiceáil mé na cealla folmha in aice láimhe: (7,6) folamh, (6,6) folamh, (5,7) folamh, (4,6) folamh. An bhfuil ceann ar bith díobh seo oiriúnach duitse?} \\
            }
        }
    }
    & & \\ \\

    \theutterance \stepcounter{utterance}  
    & & \multicolumn{4}{p{0.6\linewidth}}{
        \cellcolor[rgb]{0.9,0.9,0.9}{
            \makecell[{{p{\linewidth}}}]{
                \texttt{\tiny{[P1$\langle$GM]}}
                \texttt{Bogadh 'D' go (7,5) go rathúil. Seo mar atá do ghreille nuashonraithe:} \\
\texttt{```} \\
\texttt{ 1234567} \\
\texttt{╔═╤══╤══╗} \\
\texttt{║◌│◌◌│◌◌║ 1} \\
\texttt{║T└┬─┴──╢ 2} \\
\texttt{║◌◌│◌O◌◌║ 3} \\
\texttt{╟─┬┘SP┌─╢ 4} \\
\texttt{║◌│A◌◌│D║ 5} \\
\texttt{╟─┼─┬─┤M║ 6} \\
\texttt{║◌│◌│◌│◌║ 7} \\
\texttt{╚═╧═╧═╧═╝} \\
\texttt{```} \\
\texttt{Tá sibh ag imirt babhta 11 faoi láthair as uasmhéid 28.} \\
\texttt{Tá 10 as 16 pionós carntha agaibh le chéile.} \\
\texttt{Sheol an t{-}imreoir eile an teachtaireacht seo:} \\
\\ 
\\ 
\texttt{"Tá deacracht agam A a bhogadh. Ar mo chlár, tá A ag (7,5) faoi láthair. Tá (5,6) blocáilte ag '─', (6,5) ag '│', agus (7,4) ag '─'. An féidir linn sprioc eile a roghnú do A? Seiceáil mé na cealla folmha in aice láimhe: (7,6) folamh, (6,6) folamh, (5,7) folamh, (4,6) folamh. An bhfuil ceann ar bith díobh seo oiriúnach duitse?"} \\
\\ 
\\ 
\texttt{Cad é do chéad ordú eile?} \\
            }
        }
    }
    & \\ \\

    \theutterance \stepcounter{utterance}  
    & \multicolumn{4}{p{0.6\linewidth}}{
        \cellcolor[rgb]{0.9,0.9,0.9}{
            \makecell[{{p{\linewidth}}}]{
                \texttt{\tiny{[P1$\rangle$GM]}}
                \texttt{ABAIR: Is féidir linn A a chur ag (4,6) {-} tá sé folamh ar mo chlár. Nó (7,6) atá folamh freisin. Roghnaigh ceann acu agus cuir in iúl dom. Idir an dá linn, bogfaidh mé A ó (2,5) go dtí an sprioc a roghnaíonn muid. Tá D agam anois ag (7,5), mar sin beidh orainn sprioc D a choigeartú freisin {-} an féidir leat (5,4) nó (6,4) a úsáid do D? Nó an bhfuil do D fós ag (5,5)?} \\
            }
        }
    }
    & & \\ \\

    \theutterance \stepcounter{utterance}  
    & & & \multicolumn{2}{p{0.3\linewidth}}{
        \cellcolor[rgb]{0.9,0.9,0.9}{
            \makecell[{{p{\linewidth}}}]{
                \texttt{\tiny{[GM$|$GM]}}
                \texttt{ABAIR: Is féidir linn A a chur ag (4,6) {-} tá sé folamh ar mo chlár. Nó (7,6) atá folamh freisin. Roghnaigh ceann acu agus cuir in iúl dom. Idir an dá linn, bogfaidh mé A ó (2,5) go dtí an sprioc a roghnaíonn muid. Tá D agam anois ag (7,5), mar sin beidh orainn sprioc D a choigeartú freisin {-} an féidir leat (5,4) nó (6,4) a úsáid do D? Nó an bhfuil do D fós ag (5,5)?} \\
            }
        }
    }
    & & \\ \\

    \theutterance \stepcounter{utterance}  
    & & \multicolumn{4}{p{0.6\linewidth}}{
        \cellcolor[rgb]{0.9,0.9,0.9}{
            \makecell[{{p{\linewidth}}}]{
                \texttt{\tiny{[GM$\rangle$P2]}}
                \texttt{Cuireadh do theachtaireacht ar aghaidh chuig an imreoir eile.} \\
\\ 
\\ 
\texttt{Tá sibh ag imirt babhta 12 faoi láthair as uasmhéid 28.} \\
\texttt{Tá 10 as 16 pionós carntha agaibh le chéile.} \\
\texttt{Sheol an t{-}imreoir eile an teachtaireacht seo:} \\
\\ 
\\ 
\texttt{"Is féidir linn A a chur ag (4,6) {-} tá sé folamh ar mo chlár. Nó (7,6) atá folamh freisin. Roghnaigh ceann acu agus cuir in iúl dom. Idir an dá linn, bogfaidh mé A ó (2,5) go dtí an sprioc a roghnaíonn muid. Tá D agam anois ag (7,5), mar sin beidh orainn sprioc D a choigeartú freisin {-} an féidir leat (5,4) nó (6,4) a úsáid do D? Nó an bhfuil do D fós ag (5,5)?"} \\
\\ 
\\ 
\texttt{Cad é do chéad ordú eile?} \\
            }
        }
    }
    & \\ \\

    \theutterance \stepcounter{utterance}  
    & & & \multicolumn{4}{p{0.6\linewidth}}{
        \cellcolor[rgb]{0.9,0.9,0.9}{
            \makecell[{{p{\linewidth}}}]{
                \texttt{\tiny{[GM$\langle$P2]}}
                \texttt{ABAIR: Roghnaímis A ag (4,6). Bogfaidh mé A ann. Tá mo D ag (5,5) faoi láthair, agus is é (6,4) an sprioc a bhí againn do D ach bhí '┌' ann níos luaithe. Lig dom seiceáil arís: an bhfuil (6,4) folamh anois? Ar mo chlár, tar éis dom S a bhogadh go (4,4), tá (6,4) fós ag '┌'. Mar sin ní féidir liom D a bhogadh go (6,4). An féidir leat D a chur ag (6,4) ar do chlár? Nó ar cheart dúinn D a fhágáil ag (5,5) mar sprioc? Tá (5,5) folamh ar do chlár anois?} \\
            }
        }
    }
     \\ \\

    \theutterance \stepcounter{utterance}  
    & & & \multicolumn{2}{p{0.3\linewidth}}{
        \cellcolor[rgb]{0.9,0.9,0.9}{
            \makecell[{{p{\linewidth}}}]{
                \texttt{\tiny{[GM$|$GM]}}
                \texttt{ABAIR: Roghnaímis A ag (4,6). Bogfaidh mé A ann. Tá mo D ag (5,5) faoi láthair, agus is é (6,4) an sprioc a bhí againn do D ach bhí '┌' ann níos luaithe. Lig dom seiceáil arís: an bhfuil (6,4) folamh anois? Ar mo chlár, tar éis dom S a bhogadh go (4,4), tá (6,4) fós ag '┌'. Mar sin ní féidir liom D a bhogadh go (6,4). An féidir leat D a chur ag (6,4) ar do chlár? Nó ar cheart dúinn D a fhágáil ag (5,5) mar sprioc? Tá (5,5) folamh ar do chlár anois?} \\
            }
        }
    }
    & & \\ \\

    \theutterance \stepcounter{utterance}  
    & & \multicolumn{4}{p{0.6\linewidth}}{
        \cellcolor[rgb]{0.9,0.9,0.9}{
            \makecell[{{p{\linewidth}}}]{
                \texttt{\tiny{[P1$\langle$GM]}}
                \texttt{Cuireadh do theachtaireacht ar aghaidh chuig an imreoir eile.} \\
\\ 
\\ 
\texttt{Tá sibh ag imirt babhta 12 faoi láthair as uasmhéid 28.} \\
\texttt{Tá 10 as 16 pionós carntha agaibh le chéile.} \\
\texttt{Sheol an t{-}imreoir eile an teachtaireacht seo:} \\
\\ 
\\ 
\texttt{"Roghnaímis A ag (4,6). Bogfaidh mé A ann. Tá mo D ag (5,5) faoi láthair, agus is é (6,4) an sprioc a bhí againn do D ach bhí '┌' ann níos luaithe. Lig dom seiceáil arís: an bhfuil (6,4) folamh anois? Ar mo chlár, tar éis dom S a bhogadh go (4,4), tá (6,4) fós ag '┌'. Mar sin ní féidir liom D a bhogadh go (6,4). An féidir leat D a chur ag (6,4) ar do chlár? Nó ar cheart dúinn D a fhágáil ag (5,5) mar sprioc? Tá (5,5) folamh ar do chlár anois?"} \\
\\ 
\\ 
\texttt{Cad é do chéad ordú eile?} \\
            }
        }
    }
    & \\ \\

    \theutterance \stepcounter{utterance}  
    & \multicolumn{4}{p{0.6\linewidth}}{
        \cellcolor[rgb]{0.9,0.9,0.9}{
            \makecell[{{p{\linewidth}}}]{
                \texttt{\tiny{[P1$\rangle$GM]}}
                \texttt{ABAIR: Tá (5,5) folamh ar mo chlár anois tar éis dom D a bhogadh amach. Fágaimis D ag (5,5) mar sprioc {-} tá do D ann cheana féin, agus is féidir liom mo D a bhogadh ó (7,5) ar ais go (5,5). Bogfaidh mé A go (4,6) anois. An bhfuil T bogtha agat go (6,3) fós?} \\
            }
        }
    }
    & & \\ \\

    \theutterance \stepcounter{utterance}  
    & & & \multicolumn{2}{p{0.3\linewidth}}{
        \cellcolor[rgb]{0.9,0.9,0.9}{
            \makecell[{{p{\linewidth}}}]{
                \texttt{\tiny{[GM$|$GM]}}
                \texttt{ABAIR: Tá (5,5) folamh ar mo chlár anois tar éis dom D a bhogadh amach. Fágaimis D ag (5,5) mar sprioc {-} tá do D ann cheana féin, agus is féidir liom mo D a bhogadh ó (7,5) ar ais go (5,5). Bogfaidh mé A go (4,6) anois. An bhfuil T bogtha agat go (6,3) fós?} \\
            }
        }
    }
    & & \\ \\

    \theutterance \stepcounter{utterance}  
    & & \multicolumn{4}{p{0.6\linewidth}}{
        \cellcolor[rgb]{0.9,0.9,0.9}{
            \makecell[{{p{\linewidth}}}]{
                \texttt{\tiny{[GM$\rangle$P2]}}
                \texttt{Cuireadh do theachtaireacht ar aghaidh chuig an imreoir eile.} \\
\\ 
\\ 
\texttt{Tá sibh ag imirt babhta 13 faoi láthair as uasmhéid 28.} \\
\texttt{Tá 10 as 16 pionós carntha agaibh le chéile.} \\
\texttt{Sheol an t{-}imreoir eile an teachtaireacht seo:} \\
\\ 
\\ 
\texttt{"Tá (5,5) folamh ar mo chlár anois tar éis dom D a bhogadh amach. Fágaimis D ag (5,5) mar sprioc {-} tá do D ann cheana féin, agus is féidir liom mo D a bhogadh ó (7,5) ar ais go (5,5). Bogfaidh mé A go (4,6) anois. An bhfuil T bogtha agat go (6,3) fós?"} \\
\\ 
\\ 
\texttt{Cad é do chéad ordú eile?} \\
            }
        }
    }
    & \\ \\

    \theutterance \stepcounter{utterance}  
    & & & \multicolumn{4}{p{0.6\linewidth}}{
        \cellcolor[rgb]{0.9,0.9,0.9}{
            \makecell[{{p{\linewidth}}}]{
                \texttt{\tiny{[GM$\langle$P2]}}
                \texttt{BOG: A, (4,6)} \\
            }
        }
    }
     \\ \\

    \theutterance \stepcounter{utterance}  
    & & & \multicolumn{2}{p{0.3\linewidth}}{
        \cellcolor[rgb]{0.9,0.9,0.9}{
            \makecell[{{p{\linewidth}}}]{
                \texttt{\tiny{[GM$|$GM]}}
                \texttt{BOG: A, (4,6)} \\
            }
        }
    }
    & & \\ \\

    \theutterance \stepcounter{utterance}  
    & & & \multicolumn{2}{p{0.3\linewidth}}{
        \cellcolor[rgb]{0.9,0.9,0.9}{
            \makecell[{{p{\linewidth}}}]{
                \texttt{\tiny{[GM$|$GM]}}
                \texttt{Pionós: níl (4,6) folamh, ach tá '┬' ann.} \\
\\ 
\texttt{Tá 11 as 16 pionós carntha agaibh le chéile.} \\
\texttt{Déan iarracht arís, le do thoil!} \\
            }
        }
    }
    & & \\ \\

    \theutterance \stepcounter{utterance}  
    & & \multicolumn{4}{p{0.6\linewidth}}{
        \cellcolor[rgb]{0.9,0.9,0.9}{
            \makecell[{{p{\linewidth}}}]{
                \texttt{\tiny{[GM$\rangle$P2]}}
                \texttt{Pionós: níl (4,6) folamh, ach tá '┬' ann.} \\
\\ 
\texttt{Tá 11 as 16 pionós carntha agaibh le chéile.} \\
\texttt{Déan iarracht arís, le do thoil!} \\
            }
        }
    }
    & \\ \\

    \theutterance \stepcounter{utterance}  
    & & & \multicolumn{4}{p{0.6\linewidth}}{
        \cellcolor[rgb]{0.9,0.9,0.9}{
            \makecell[{{p{\linewidth}}}]{
                \texttt{\tiny{[GM$\langle$P2]}}
                \texttt{BOG: A, (7,6)} \\
            }
        }
    }
     \\ \\

    \theutterance \stepcounter{utterance}  
    & & & \multicolumn{2}{p{0.3\linewidth}}{
        \cellcolor[rgb]{0.9,0.9,0.9}{
            \makecell[{{p{\linewidth}}}]{
                \texttt{\tiny{[GM$|$GM]}}
                \texttt{BOG: A, (7,6)} \\
            }
        }
    }
    & & \\ \\

    \theutterance \stepcounter{utterance}  
    & & & \multicolumn{2}{p{0.3\linewidth}}{
        \cellcolor[rgb]{0.9,0.9,0.9}{
            \makecell[{{p{\linewidth}}}]{
                \texttt{\tiny{[GM$|$GM]}}
                \texttt{A moved to (7, 6)} \\
            }
        }
    }
    & & \\ \\

    \theutterance \stepcounter{utterance}  
    & & \multicolumn{4}{p{0.6\linewidth}}{
        \cellcolor[rgb]{0.9,0.9,0.9}{
            \makecell[{{p{\linewidth}}}]{
                \texttt{\tiny{[P1$\langle$GM]}}
                \texttt{Cuireadh do theachtaireacht ar aghaidh chuig an imreoir eile.} \\
\\ 
\\ 
\texttt{Tá sibh ag imirt babhta 13 faoi láthair as uasmhéid 28.} \\
\texttt{Tá 11 as 16 pionós carntha agaibh le chéile.} \\
\texttt{Bhog an t{-}imreoir eile réad ar a ngreille.} \\
\\ 
\\ 
\texttt{Cad é do chéad ordú eile?} \\
            }
        }
    }
    & \\ \\

    \theutterance \stepcounter{utterance}  
    & \multicolumn{4}{p{0.6\linewidth}}{
        \cellcolor[rgb]{0.9,0.9,0.9}{
            \makecell[{{p{\linewidth}}}]{
                \texttt{\tiny{[P1$\rangle$GM]}}
                \texttt{BOG: A, (4,6)} \\
            }
        }
    }
    & & \\ \\

    \theutterance \stepcounter{utterance}  
    & & & \multicolumn{2}{p{0.3\linewidth}}{
        \cellcolor[rgb]{0.9,0.9,0.9}{
            \makecell[{{p{\linewidth}}}]{
                \texttt{\tiny{[GM$|$GM]}}
                \texttt{BOG: A, (4,6)} \\
            }
        }
    }
    & & \\ \\

    \theutterance \stepcounter{utterance}  
    & & & \multicolumn{2}{p{0.3\linewidth}}{
        \cellcolor[rgb]{0.9,0.9,0.9}{
            \makecell[{{p{\linewidth}}}]{
                \texttt{\tiny{[GM$|$GM]}}
                \texttt{Pionós: níl (4,6) folamh, ach tá '┬' ann.} \\
\\ 
\texttt{Tá 12 as 16 pionós carntha agaibh le chéile.} \\
\texttt{Déan iarracht arís, le do thoil!} \\
            }
        }
    }
    & & \\ \\

    \theutterance \stepcounter{utterance}  
    & & \multicolumn{4}{p{0.6\linewidth}}{
        \cellcolor[rgb]{0.9,0.9,0.9}{
            \makecell[{{p{\linewidth}}}]{
                \texttt{\tiny{[P1$\langle$GM]}}
                \texttt{Pionós: níl (4,6) folamh, ach tá '┬' ann.} \\
\\ 
\texttt{Tá 12 as 16 pionós carntha agaibh le chéile.} \\
\texttt{Déan iarracht arís, le do thoil!} \\
            }
        }
    }
    & \\ \\

    \theutterance \stepcounter{utterance}  
    & \multicolumn{4}{p{0.6\linewidth}}{
        \cellcolor[rgb]{0.9,0.9,0.9}{
            \makecell[{{p{\linewidth}}}]{
                \texttt{\tiny{[P1$\rangle$GM]}}
                \texttt{BOG: A, (7,6)} \\
            }
        }
    }
    & & \\ \\

    \theutterance \stepcounter{utterance}  
    & & & \multicolumn{2}{p{0.3\linewidth}}{
        \cellcolor[rgb]{0.9,0.9,0.9}{
            \makecell[{{p{\linewidth}}}]{
                \texttt{\tiny{[GM$|$GM]}}
                \texttt{BOG: A, (7,6)} \\
            }
        }
    }
    & & \\ \\

    \theutterance \stepcounter{utterance}  
    & & & \multicolumn{2}{p{0.3\linewidth}}{
        \cellcolor[rgb]{0.9,0.9,0.9}{
            \makecell[{{p{\linewidth}}}]{
                \texttt{\tiny{[GM$|$GM]}}
                \texttt{Pionós: níl (7,6) folamh, ach tá 'M' ann.} \\
\\ 
\texttt{Tá 13 as 16 pionós carntha agaibh le chéile.} \\
\texttt{Déan iarracht arís, le do thoil!} \\
            }
        }
    }
    & & \\ \\

    \theutterance \stepcounter{utterance}  
    & & \multicolumn{4}{p{0.6\linewidth}}{
        \cellcolor[rgb]{0.9,0.9,0.9}{
            \makecell[{{p{\linewidth}}}]{
                \texttt{\tiny{[P1$\langle$GM]}}
                \texttt{Pionós: níl (7,6) folamh, ach tá 'M' ann.} \\
\\ 
\texttt{Tá 13 as 16 pionós carntha agaibh le chéile.} \\
\texttt{Déan iarracht arís, le do thoil!} \\
            }
        }
    }
    & \\ \\

    \theutterance \stepcounter{utterance}  
    & \multicolumn{4}{p{0.6\linewidth}}{
        \cellcolor[rgb]{0.9,0.9,0.9}{
            \makecell[{{p{\linewidth}}}]{
                \texttt{\tiny{[P1$\rangle$GM]}}
                \texttt{BOG: A, (1,1)} \\
            }
        }
    }
    & & \\ \\

    \theutterance \stepcounter{utterance}  
    & & & \multicolumn{2}{p{0.3\linewidth}}{
        \cellcolor[rgb]{0.9,0.9,0.9}{
            \makecell[{{p{\linewidth}}}]{
                \texttt{\tiny{[GM$|$GM]}}
                \texttt{BOG: A, (1,1)} \\
            }
        }
    }
    & & \\ \\

    \theutterance \stepcounter{utterance}  
    & & & \multicolumn{2}{p{0.3\linewidth}}{
        \cellcolor[rgb]{0.9,0.9,0.9}{
            \makecell[{{p{\linewidth}}}]{
                \texttt{\tiny{[GM$|$GM]}}
                \texttt{A moved to (1, 1)} \\
            }
        }
    }
    & & \\ \\

    \theutterance \stepcounter{utterance}  
    & & \multicolumn{4}{p{0.6\linewidth}}{
        \cellcolor[rgb]{0.9,0.9,0.9}{
            \makecell[{{p{\linewidth}}}]{
                \texttt{\tiny{[GM$\rangle$P2]}}
                \texttt{Bogadh 'A' go (7,6) go rathúil. Seo mar atá do ghreille nuashonraithe:} \\
\texttt{```} \\
\texttt{ 1234567} \\
\texttt{╔═╤══╤══╗} \\
\texttt{║◌│◌◌│◌◌║ 1} \\
\texttt{║◌└┬─┴──╢ 2} \\
\texttt{║◌◌│◌OT◌║ 3} \\
\texttt{╟─┬┘SP┌─╢ 4} \\
\texttt{║◌│◌◌D│◌║ 5} \\
\texttt{╟─┼─┬─┤A║ 6} \\
\texttt{║M│◌│◌│◌║ 7} \\
\texttt{╚═╧═╧═╧═╝} \\
\texttt{```} \\
\texttt{Tá sibh ag imirt babhta 14 faoi láthair as uasmhéid 28.} \\
\texttt{Tá 13 as 16 pionós carntha agaibh le chéile.} \\
\texttt{Bhog an t{-}imreoir eile réad ar a ngreille.} \\
\\ 
\\ 
\texttt{Cad é do chéad ordú eile?} \\
            }
        }
    }
    & \\ \\

    \theutterance \stepcounter{utterance}  
    & & & \multicolumn{4}{p{0.6\linewidth}}{
        \cellcolor[rgb]{0.9,0.9,0.9}{
            \makecell[{{p{\linewidth}}}]{
                \texttt{\tiny{[GM$\langle$P2]}}
                \texttt{ABAIR: D'éirigh liom A a bhogadh go (7,6). Tá T ag (6,3) cheana féin. Anois tá P ag (5,4), O ag (5,3), T ag (6,3), S ag (4,4), D ag (5,5), A ag (7,6), agus M ag (1,7). Cad iad na suímh atá agat anois? Agus cad faoi M {-} cad é an sprioc do M? Ar cheart M a bhogadh go (5,7) nó (6,7) nó áit éigin eile?} \\
            }
        }
    }
     \\ \\

    \theutterance \stepcounter{utterance}  
    & & & \multicolumn{2}{p{0.3\linewidth}}{
        \cellcolor[rgb]{0.9,0.9,0.9}{
            \makecell[{{p{\linewidth}}}]{
                \texttt{\tiny{[GM$|$GM]}}
                \texttt{ABAIR: D'éirigh liom A a bhogadh go (7,6). Tá T ag (6,3) cheana féin. Anois tá P ag (5,4), O ag (5,3), T ag (6,3), S ag (4,4), D ag (5,5), A ag (7,6), agus M ag (1,7). Cad iad na suímh atá agat anois? Agus cad faoi M {-} cad é an sprioc do M? Ar cheart M a bhogadh go (5,7) nó (6,7) nó áit éigin eile?} \\
            }
        }
    }
    & & \\ \\

    \theutterance \stepcounter{utterance}  
    & & \multicolumn{4}{p{0.6\linewidth}}{
        \cellcolor[rgb]{0.9,0.9,0.9}{
            \makecell[{{p{\linewidth}}}]{
                \texttt{\tiny{[P1$\langle$GM]}}
                \texttt{Bogadh 'A' go (1,1) go rathúil. Seo mar atá do ghreille nuashonraithe:} \\
\texttt{```} \\
\texttt{ 1234567} \\
\texttt{╔═╤══╤══╗} \\
\texttt{║A│◌◌│◌◌║ 1} \\
\texttt{║T└┬─┴──╢ 2} \\
\texttt{║◌◌│◌O◌◌║ 3} \\
\texttt{╟─┬┘SP┌─╢ 4} \\
\texttt{║◌│◌◌◌│D║ 5} \\
\texttt{╟─┼─┬─┤M║ 6} \\
\texttt{║◌│◌│◌│◌║ 7} \\
\texttt{╚═╧═╧═╧═╝} \\
\texttt{```} \\
\texttt{Tá sibh ag imirt babhta 14 faoi láthair as uasmhéid 28.} \\
\texttt{Tá 13 as 16 pionós carntha agaibh le chéile.} \\
\texttt{Sheol an t{-}imreoir eile an teachtaireacht seo:} \\
\\ 
\\ 
\texttt{"D'éirigh liom A a bhogadh go (7,6). Tá T ag (6,3) cheana féin. Anois tá P ag (5,4), O ag (5,3), T ag (6,3), S ag (4,4), D ag (5,5), A ag (7,6), agus M ag (1,7). Cad iad na suímh atá agat anois? Agus cad faoi M {-} cad é an sprioc do M? Ar cheart M a bhogadh go (5,7) nó (6,7) nó áit éigin eile?"} \\
\\ 
\\ 
\texttt{Cad é do chéad ordú eile?} \\
            }
        }
    }
    & \\ \\

    \theutterance \stepcounter{utterance}  
    & \multicolumn{4}{p{0.6\linewidth}}{
        \cellcolor[rgb]{0.9,0.9,0.9}{
            \makecell[{{p{\linewidth}}}]{
                \texttt{\tiny{[P1$\rangle$GM]}}
                \texttt{ABAIR: Tá P(5,4), O(5,3), S(4,4), D(7,5), M(7,6) agam. Tá T ag (1,2) agus A ag (1,1) faoi láthair. Ní féidir liom A a chur ag (4,6) mar gheall ar '┬', agus (7,6) tá M ann. An féidir linn A a chur ag (1,4) nó (2,6) ina ionad? Seiceáil do chlár. Do M, roghnaímis (5,6) nó (6,5) {-} an bhfuil ceann acu folamh agat? Ní mór dúinn a bheith cúramach {-} níl ach 3 phionós fágtha againn!} \\
            }
        }
    }
    & & \\ \\

    \theutterance \stepcounter{utterance}  
    & & & \multicolumn{2}{p{0.3\linewidth}}{
        \cellcolor[rgb]{0.9,0.9,0.9}{
            \makecell[{{p{\linewidth}}}]{
                \texttt{\tiny{[GM$|$GM]}}
                \texttt{ABAIR: Tá P(5,4), O(5,3), S(4,4), D(7,5), M(7,6) agam. Tá T ag (1,2) agus A ag (1,1) faoi láthair. Ní féidir liom A a chur ag (4,6) mar gheall ar '┬', agus (7,6) tá M ann. An féidir linn A a chur ag (1,4) nó (2,6) ina ionad? Seiceáil do chlár. Do M, roghnaímis (5,6) nó (6,5) {-} an bhfuil ceann acu folamh agat? Ní mór dúinn a bheith cúramach {-} níl ach 3 phionós fágtha againn!} \\
            }
        }
    }
    & & \\ \\

    \theutterance \stepcounter{utterance}  
    & & \multicolumn{4}{p{0.6\linewidth}}{
        \cellcolor[rgb]{0.9,0.9,0.9}{
            \makecell[{{p{\linewidth}}}]{
                \texttt{\tiny{[GM$\rangle$P2]}}
                \texttt{Cuireadh do theachtaireacht ar aghaidh chuig an imreoir eile.} \\
\\ 
\\ 
\texttt{Tá sibh ag imirt babhta 15 faoi láthair as uasmhéid 28.} \\
\texttt{Tá 13 as 16 pionós carntha agaibh le chéile.} \\
\texttt{Sheol an t{-}imreoir eile an teachtaireacht seo:} \\
\\ 
\\ 
\texttt{"Tá P(5,4), O(5,3), S(4,4), D(7,5), M(7,6) agam. Tá T ag (1,2) agus A ag (1,1) faoi láthair. Ní féidir liom A a chur ag (4,6) mar gheall ar '┬', agus (7,6) tá M ann. An féidir linn A a chur ag (1,4) nó (2,6) ina ionad? Seiceáil do chlár. Do M, roghnaímis (5,6) nó (6,5) {-} an bhfuil ceann acu folamh agat? Ní mór dúinn a bheith cúramach {-} níl ach 3 phionós fágtha againn!"} \\
\\ 
\\ 
\texttt{Cad é do chéad ordú eile?} \\
            }
        }
    }
    & \\ \\

    \theutterance \stepcounter{utterance}  
    & & & \multicolumn{4}{p{0.6\linewidth}}{
        \cellcolor[rgb]{0.9,0.9,0.9}{
            \makecell[{{p{\linewidth}}}]{
                \texttt{\tiny{[GM$\langle$P2]}}
                \texttt{ABAIR: Seiceáil mé mo chlár: (1,4) tá '┘' ann, (2,6) tá '┼' ann, mar sin níl ceachtar acu folamh. Maidir le A, cad faoi (2,2) nó (3,2) nó (7,2)? Ar mo chlár: (2,2) tá '└' ann, (3,2) tá '┬' ann, (7,2) folamh. Do M: (5,6) tá '┤' ann, (6,5) tá '│' ann. Cad faoi M ag (6,7) nó (7,7)? Tá (6,7) folamh agus (7,7) folamh ar mo chlár. Nó b'fhéidir A ag (1,3)? Tá (1,3) folamh. Caithfimid na spriocanna a choigeartú go cúramach!} \\
            }
        }
    }
     \\ \\

    \theutterance \stepcounter{utterance}  
    & & & \multicolumn{2}{p{0.3\linewidth}}{
        \cellcolor[rgb]{0.9,0.9,0.9}{
            \makecell[{{p{\linewidth}}}]{
                \texttt{\tiny{[GM$|$GM]}}
                \texttt{ABAIR: Seiceáil mé mo chlár: (1,4) tá '┘' ann, (2,6) tá '┼' ann, mar sin níl ceachtar acu folamh. Maidir le A, cad faoi (2,2) nó (3,2) nó (7,2)? Ar mo chlár: (2,2) tá '└' ann, (3,2) tá '┬' ann, (7,2) folamh. Do M: (5,6) tá '┤' ann, (6,5) tá '│' ann. Cad faoi M ag (6,7) nó (7,7)? Tá (6,7) folamh agus (7,7) folamh ar mo chlár. Nó b'fhéidir A ag (1,3)? Tá (1,3) folamh. Caithfimid na spriocanna a choigeartú go cúramach!} \\
            }
        }
    }
    & & \\ \\

    \theutterance \stepcounter{utterance}  
    & & \multicolumn{4}{p{0.6\linewidth}}{
        \cellcolor[rgb]{0.9,0.9,0.9}{
            \makecell[{{p{\linewidth}}}]{
                \texttt{\tiny{[P1$\langle$GM]}}
                \texttt{Cuireadh do theachtaireacht ar aghaidh chuig an imreoir eile.} \\
\\ 
\\ 
\texttt{Tá sibh ag imirt babhta 15 faoi láthair as uasmhéid 28.} \\
\texttt{Tá 13 as 16 pionós carntha agaibh le chéile.} \\
\texttt{Sheol an t{-}imreoir eile an teachtaireacht seo:} \\
\\ 
\\ 
\texttt{"Seiceáil mé mo chlár: (1,4) tá '┘' ann, (2,6) tá '┼' ann, mar sin níl ceachtar acu folamh. Maidir le A, cad faoi (2,2) nó (3,2) nó (7,2)? Ar mo chlár: (2,2) tá '└' ann, (3,2) tá '┬' ann, (7,2) folamh. Do M: (5,6) tá '┤' ann, (6,5) tá '│' ann. Cad faoi M ag (6,7) nó (7,7)? Tá (6,7) folamh agus (7,7) folamh ar mo chlár. Nó b'fhéidir A ag (1,3)? Tá (1,3) folamh. Caithfimid na spriocanna a choigeartú go cúramach!"} \\
\\ 
\\ 
\texttt{Cad é do chéad ordú eile?} \\
            }
        }
    }
    & \\ \\

    \theutterance \stepcounter{utterance}  
    & \multicolumn{4}{p{0.6\linewidth}}{
        \cellcolor[rgb]{0.9,0.9,0.9}{
            \makecell[{{p{\linewidth}}}]{
                \texttt{\tiny{[P1$\rangle$GM]}}
                \texttt{ABAIR: Roghnaímis A ag (1,3) agus M ag (6,7). Tá (1,3) folamh ar mo chlár, agus (6,7) folamh freisin. Bogfaidh mé A ó (1,1) go (1,3) agus M ó (7,6) go (6,7). Ansin caithfidh mé D a bhogadh ó (7,5) go (5,5) agus T ó (1,2) go (6,3). An bhfuil na spriocanna seo go léir folamh ar do chlár: A(1,3), M(6,7), D(5,5), T(6,3)?} \\
            }
        }
    }
    & & \\ \\

    \theutterance \stepcounter{utterance}  
    & & & \multicolumn{2}{p{0.3\linewidth}}{
        \cellcolor[rgb]{0.9,0.9,0.9}{
            \makecell[{{p{\linewidth}}}]{
                \texttt{\tiny{[GM$|$GM]}}
                \texttt{ABAIR: Roghnaímis A ag (1,3) agus M ag (6,7). Tá (1,3) folamh ar mo chlár, agus (6,7) folamh freisin. Bogfaidh mé A ó (1,1) go (1,3) agus M ó (7,6) go (6,7). Ansin caithfidh mé D a bhogadh ó (7,5) go (5,5) agus T ó (1,2) go (6,3). An bhfuil na spriocanna seo go léir folamh ar do chlár: A(1,3), M(6,7), D(5,5), T(6,3)?} \\
            }
        }
    }
    & & \\ \\

    \theutterance \stepcounter{utterance}  
    & & \multicolumn{4}{p{0.6\linewidth}}{
        \cellcolor[rgb]{0.9,0.9,0.9}{
            \makecell[{{p{\linewidth}}}]{
                \texttt{\tiny{[GM$\rangle$P2]}}
                \texttt{Cuireadh do theachtaireacht ar aghaidh chuig an imreoir eile.} \\
\\ 
\\ 
\texttt{Tá sibh ag imirt babhta 16 faoi láthair as uasmhéid 28.} \\
\texttt{Tá 13 as 16 pionós carntha agaibh le chéile.} \\
\texttt{Sheol an t{-}imreoir eile an teachtaireacht seo:} \\
\\ 
\\ 
\texttt{"Roghnaímis A ag (1,3) agus M ag (6,7). Tá (1,3) folamh ar mo chlár, agus (6,7) folamh freisin. Bogfaidh mé A ó (1,1) go (1,3) agus M ó (7,6) go (6,7). Ansin caithfidh mé D a bhogadh ó (7,5) go (5,5) agus T ó (1,2) go (6,3). An bhfuil na spriocanna seo go léir folamh ar do chlár: A(1,3), M(6,7), D(5,5), T(6,3)?"} \\
\\ 
\\ 
\texttt{Cad é do chéad ordú eile?} \\
            }
        }
    }
    & \\ \\

    \theutterance \stepcounter{utterance}  
    & & & \multicolumn{4}{p{0.6\linewidth}}{
        \cellcolor[rgb]{0.9,0.9,0.9}{
            \makecell[{{p{\linewidth}}}]{
                \texttt{\tiny{[GM$\langle$P2]}}
                \texttt{ABAIR: Deimhním: Tá T agam cheana ag (6,3) agus D ag (5,5). Tá (1,3) folamh agus (6,7) folamh ar mo chlár. Mar sin tá na spriocanna go léir oiriúnach! Níl fágtha agam ach M a bhogadh ó (1,7) go (6,7). Nuair a bheidh sin déanta, beidh mo rudaí go léir ag na sprioc{-}shuímh: P(5,4), O(5,3), T(6,3), S(4,4), D(5,5), A(7,6)... fan nóiméad, tá A agam ag (7,6) cheana ach is é (1,3) an sprioc nua! Caithfidh mé A a bhogadh ó (7,6) go (1,3) freisin. An bhfuil (7,6) folamh ar do chlár le haghaidh do M? Agus an bhfuil (1,3) folamh duitse le haghaidh A?} \\
            }
        }
    }
     \\ \\

    \theutterance \stepcounter{utterance}  
    & & & \multicolumn{2}{p{0.3\linewidth}}{
        \cellcolor[rgb]{0.9,0.9,0.9}{
            \makecell[{{p{\linewidth}}}]{
                \texttt{\tiny{[GM$|$GM]}}
                \texttt{ABAIR: Deimhním: Tá T agam cheana ag (6,3) agus D ag (5,5). Tá (1,3) folamh agus (6,7) folamh ar mo chlár. Mar sin tá na spriocanna go léir oiriúnach! Níl fágtha agam ach M a bhogadh ó (1,7) go (6,7). Nuair a bheidh sin déanta, beidh mo rudaí go léir ag na sprioc{-}shuímh: P(5,4), O(5,3), T(6,3), S(4,4), D(5,5), A(7,6)... fan nóiméad, tá A agam ag (7,6) cheana ach is é (1,3) an sprioc nua! Caithfidh mé A a bhogadh ó (7,6) go (1,3) freisin. An bhfuil (7,6) folamh ar do chlár le haghaidh do M? Agus an bhfuil (1,3) folamh duitse le haghaidh A?} \\
            }
        }
    }
    & & \\ \\

    \theutterance \stepcounter{utterance}  
    & & \multicolumn{4}{p{0.6\linewidth}}{
        \cellcolor[rgb]{0.9,0.9,0.9}{
            \makecell[{{p{\linewidth}}}]{
                \texttt{\tiny{[P1$\langle$GM]}}
                \texttt{Cuireadh do theachtaireacht ar aghaidh chuig an imreoir eile.} \\
\\ 
\\ 
\texttt{Tá sibh ag imirt babhta 16 faoi láthair as uasmhéid 28.} \\
\texttt{Tá 13 as 16 pionós carntha agaibh le chéile.} \\
\texttt{Sheol an t{-}imreoir eile an teachtaireacht seo:} \\
\\ 
\\ 
\texttt{"Deimhním: Tá T agam cheana ag (6,3) agus D ag (5,5). Tá (1,3) folamh agus (6,7) folamh ar mo chlár. Mar sin tá na spriocanna go léir oiriúnach! Níl fágtha agam ach M a bhogadh ó (1,7) go (6,7). Nuair a bheidh sin déanta, beidh mo rudaí go léir ag na sprioc{-}shuímh: P(5,4), O(5,3), T(6,3), S(4,4), D(5,5), A(7,6)... fan nóiméad, tá A agam ag (7,6) cheana ach is é (1,3) an sprioc nua! Caithfidh mé A a bhogadh ó (7,6) go (1,3) freisin. An bhfuil (7,6) folamh ar do chlár le haghaidh do M? Agus an bhfuil (1,3) folamh duitse le haghaidh A?"} \\
\\ 
\\ 
\texttt{Cad é do chéad ordú eile?} \\
            }
        }
    }
    & \\ \\

    \theutterance \stepcounter{utterance}  
    & \multicolumn{4}{p{0.6\linewidth}}{
        \cellcolor[rgb]{0.9,0.9,0.9}{
            \makecell[{{p{\linewidth}}}]{
                \texttt{\tiny{[P1$\rangle$GM]}}
                \texttt{ABAIR: Fan nóiméad {-} déanaimis soiléiriú. Is iad na sprioc{-}shuímh deiridh: P(5,4), O(5,3), T(6,3), S(4,4), D(5,5), A(1,3), M(6,7). Tá (7,6) folamh ar mo chlár anois tar éis dom A a bhogadh amach, agus tá (1,3) folamh. Ach má bhogann tú A go (1,3), ansin beidh (7,6) folamh agat do M? Fan {-} dúirt tú go bhfuil A agat ag (7,6) faoi láthair. Má bhogann tú A go (1,3), beidh (7,6) folamh. Ach beidh M uaitse ag dul go (6,7), ní go (7,6). Mar sin bog do A go (1,3), agus bog do M ó (1,7) go (6,7). Bogfaidh mé mo chuid rudaí de réir an phlean. Níl ach 3 phionós fágtha againn, mar sin bí cúramach!} \\
            }
        }
    }
    & & \\ \\

    \theutterance \stepcounter{utterance}  
    & & & \multicolumn{2}{p{0.3\linewidth}}{
        \cellcolor[rgb]{0.9,0.9,0.9}{
            \makecell[{{p{\linewidth}}}]{
                \texttt{\tiny{[GM$|$GM]}}
                \texttt{ABAIR: Fan nóiméad {-} déanaimis soiléiriú. Is iad na sprioc{-}shuímh deiridh: P(5,4), O(5,3), T(6,3), S(4,4), D(5,5), A(1,3), M(6,7). Tá (7,6) folamh ar mo chlár anois tar éis dom A a bhogadh amach, agus tá (1,3) folamh. Ach má bhogann tú A go (1,3), ansin beidh (7,6) folamh agat do M? Fan {-} dúirt tú go bhfuil A agat ag (7,6) faoi láthair. Má bhogann tú A go (1,3), beidh (7,6) folamh. Ach beidh M uaitse ag dul go (6,7), ní go (7,6). Mar sin bog do A go (1,3), agus bog do M ó (1,7) go (6,7). Bogfaidh mé mo chuid rudaí de réir an phlean. Níl ach 3 phionós fágtha againn, mar sin bí cúramach!} \\
            }
        }
    }
    & & \\ \\

    \theutterance \stepcounter{utterance}  
    & & \multicolumn{4}{p{0.6\linewidth}}{
        \cellcolor[rgb]{0.9,0.9,0.9}{
            \makecell[{{p{\linewidth}}}]{
                \texttt{\tiny{[GM$\rangle$P2]}}
                \texttt{Cuireadh do theachtaireacht ar aghaidh chuig an imreoir eile.} \\
\\ 
\\ 
\texttt{Tá sibh ag imirt babhta 17 faoi láthair as uasmhéid 28.} \\
\texttt{Tá 13 as 16 pionós carntha agaibh le chéile.} \\
\texttt{Sheol an t{-}imreoir eile an teachtaireacht seo:} \\
\\ 
\\ 
\texttt{"Fan nóiméad {-} déanaimis soiléiriú. Is iad na sprioc{-}shuímh deiridh: P(5,4), O(5,3), T(6,3), S(4,4), D(5,5), A(1,3), M(6,7). Tá (7,6) folamh ar mo chlár anois tar éis dom A a bhogadh amach, agus tá (1,3) folamh. Ach má bhogann tú A go (1,3), ansin beidh (7,6) folamh agat do M? Fan {-} dúirt tú go bhfuil A agat ag (7,6) faoi láthair. Má bhogann tú A go (1,3), beidh (7,6) folamh. Ach beidh M uaitse ag dul go (6,7), ní go (7,6). Mar sin bog do A go (1,3), agus bog do M ó (1,7) go (6,7). Bogfaidh mé mo chuid rudaí de réir an phlean. Níl ach 3 phionós fágtha againn, mar sin bí cúramach!"} \\
\\ 
\\ 
\texttt{Cad é do chéad ordú eile?} \\
            }
        }
    }
    & \\ \\

    \theutterance \stepcounter{utterance}  
    & & & \multicolumn{4}{p{0.6\linewidth}}{
        \cellcolor[rgb]{0.9,0.9,0.9}{
            \makecell[{{p{\linewidth}}}]{
                \texttt{\tiny{[GM$\langle$P2]}}
                \texttt{BOG: A, (1,3)} \\
            }
        }
    }
     \\ \\

    \theutterance \stepcounter{utterance}  
    & & & \multicolumn{2}{p{0.3\linewidth}}{
        \cellcolor[rgb]{0.9,0.9,0.9}{
            \makecell[{{p{\linewidth}}}]{
                \texttt{\tiny{[GM$|$GM]}}
                \texttt{BOG: A, (1,3)} \\
            }
        }
    }
    & & \\ \\

    \theutterance \stepcounter{utterance}  
    & & & \multicolumn{2}{p{0.3\linewidth}}{
        \cellcolor[rgb]{0.9,0.9,0.9}{
            \makecell[{{p{\linewidth}}}]{
                \texttt{\tiny{[GM$|$GM]}}
                \texttt{A moved to (1, 3)} \\
            }
        }
    }
    & & \\ \\

    \theutterance \stepcounter{utterance}  
    & & \multicolumn{4}{p{0.6\linewidth}}{
        \cellcolor[rgb]{0.9,0.9,0.9}{
            \makecell[{{p{\linewidth}}}]{
                \texttt{\tiny{[P1$\langle$GM]}}
                \texttt{Cuireadh do theachtaireacht ar aghaidh chuig an imreoir eile.} \\
\\ 
\\ 
\texttt{Tá sibh ag imirt babhta 17 faoi láthair as uasmhéid 28.} \\
\texttt{Tá 13 as 16 pionós carntha agaibh le chéile.} \\
\texttt{Bhog an t{-}imreoir eile réad ar a ngreille.} \\
\\ 
\\ 
\texttt{Cad é do chéad ordú eile?} \\
            }
        }
    }
    & \\ \\

    \theutterance \stepcounter{utterance}  
    & \multicolumn{4}{p{0.6\linewidth}}{
        \cellcolor[rgb]{0.9,0.9,0.9}{
            \makecell[{{p{\linewidth}}}]{
                \texttt{\tiny{[P1$\rangle$GM]}}
                \texttt{BOG: A, (1,3)} \\
            }
        }
    }
    & & \\ \\

    \theutterance \stepcounter{utterance}  
    & & & \multicolumn{2}{p{0.3\linewidth}}{
        \cellcolor[rgb]{0.9,0.9,0.9}{
            \makecell[{{p{\linewidth}}}]{
                \texttt{\tiny{[GM$|$GM]}}
                \texttt{BOG: A, (1,3)} \\
            }
        }
    }
    & & \\ \\

    \theutterance \stepcounter{utterance}  
    & & & \multicolumn{2}{p{0.3\linewidth}}{
        \cellcolor[rgb]{0.9,0.9,0.9}{
            \makecell[{{p{\linewidth}}}]{
                \texttt{\tiny{[GM$|$GM]}}
                \texttt{A moved to (1, 3)} \\
            }
        }
    }
    & & \\ \\

    \theutterance \stepcounter{utterance}  
    & & \multicolumn{4}{p{0.6\linewidth}}{
        \cellcolor[rgb]{0.9,0.9,0.9}{
            \makecell[{{p{\linewidth}}}]{
                \texttt{\tiny{[GM$\rangle$P2]}}
                \texttt{Bogadh 'A' go (1,3) go rathúil. Seo mar atá do ghreille nuashonraithe:} \\
\texttt{```} \\
\texttt{ 1234567} \\
\texttt{╔═╤══╤══╗} \\
\texttt{║◌│◌◌│◌◌║ 1} \\
\texttt{║◌└┬─┴──╢ 2} \\
\texttt{║A◌│◌OT◌║ 3} \\
\texttt{╟─┬┘SP┌─╢ 4} \\
\texttt{║◌│◌◌D│◌║ 5} \\
\texttt{╟─┼─┬─┤◌║ 6} \\
\texttt{║M│◌│◌│◌║ 7} \\
\texttt{╚═╧═╧═╧═╝} \\
\texttt{```} \\
\texttt{Tá sibh ag imirt babhta 18 faoi láthair as uasmhéid 28.} \\
\texttt{Tá 13 as 16 pionós carntha agaibh le chéile.} \\
\texttt{Bhog an t{-}imreoir eile réad ar a ngreille.} \\
\\ 
\\ 
\texttt{Cad é do chéad ordú eile?} \\
            }
        }
    }
    & \\ \\

    \theutterance \stepcounter{utterance}  
    & & & \multicolumn{4}{p{0.6\linewidth}}{
        \cellcolor[rgb]{0.9,0.9,0.9}{
            \makecell[{{p{\linewidth}}}]{
                \texttt{\tiny{[GM$\langle$P2]}}
                \texttt{BOG: M, (6,7)} \\
            }
        }
    }
     \\ \\

    \theutterance \stepcounter{utterance}  
    & & & \multicolumn{2}{p{0.3\linewidth}}{
        \cellcolor[rgb]{0.9,0.9,0.9}{
            \makecell[{{p{\linewidth}}}]{
                \texttt{\tiny{[GM$|$GM]}}
                \texttt{BOG: M, (6,7)} \\
            }
        }
    }
    & & \\ \\

    \theutterance \stepcounter{utterance}  
    & & & \multicolumn{2}{p{0.3\linewidth}}{
        \cellcolor[rgb]{0.9,0.9,0.9}{
            \makecell[{{p{\linewidth}}}]{
                \texttt{\tiny{[GM$|$GM]}}
                \texttt{Pionós: níl (6,7) folamh, ach tá '│' ann.} \\
\\ 
\texttt{Tá 14 as 16 pionós carntha agaibh le chéile.} \\
\texttt{Déan iarracht arís, le do thoil!} \\
            }
        }
    }
    & & \\ \\

    \theutterance \stepcounter{utterance}  
    & & \multicolumn{4}{p{0.6\linewidth}}{
        \cellcolor[rgb]{0.9,0.9,0.9}{
            \makecell[{{p{\linewidth}}}]{
                \texttt{\tiny{[GM$\rangle$P2]}}
                \texttt{Pionós: níl (6,7) folamh, ach tá '│' ann.} \\
\\ 
\texttt{Tá 14 as 16 pionós carntha agaibh le chéile.} \\
\texttt{Déan iarracht arís, le do thoil!} \\
            }
        }
    }
    & \\ \\

    \theutterance \stepcounter{utterance}  
    & & & \multicolumn{4}{p{0.6\linewidth}}{
        \cellcolor[rgb]{0.9,0.9,0.9}{
            \makecell[{{p{\linewidth}}}]{
                \texttt{\tiny{[GM$\langle$P2]}}
                \texttt{BOG: M, (5,7)} \\
            }
        }
    }
     \\ \\

    \theutterance \stepcounter{utterance}  
    & & & \multicolumn{2}{p{0.3\linewidth}}{
        \cellcolor[rgb]{0.9,0.9,0.9}{
            \makecell[{{p{\linewidth}}}]{
                \texttt{\tiny{[GM$|$GM]}}
                \texttt{BOG: M, (5,7)} \\
            }
        }
    }
    & & \\ \\

    \theutterance \stepcounter{utterance}  
    & & & \multicolumn{2}{p{0.3\linewidth}}{
        \cellcolor[rgb]{0.9,0.9,0.9}{
            \makecell[{{p{\linewidth}}}]{
                \texttt{\tiny{[GM$|$GM]}}
                \texttt{M moved to (5, 7)} \\
            }
        }
    }
    & & \\ \\

    \theutterance \stepcounter{utterance}  
    & & \multicolumn{4}{p{0.6\linewidth}}{
        \cellcolor[rgb]{0.9,0.9,0.9}{
            \makecell[{{p{\linewidth}}}]{
                \texttt{\tiny{[P1$\langle$GM]}}
                \texttt{Bogadh 'A' go (1,3) go rathúil. Seo mar atá do ghreille nuashonraithe:} \\
\texttt{```} \\
\texttt{ 1234567} \\
\texttt{╔═╤══╤══╗} \\
\texttt{║◌│◌◌│◌◌║ 1} \\
\texttt{║T└┬─┴──╢ 2} \\
\texttt{║A◌│◌O◌◌║ 3} \\
\texttt{╟─┬┘SP┌─╢ 4} \\
\texttt{║◌│◌◌◌│D║ 5} \\
\texttt{╟─┼─┬─┤M║ 6} \\
\texttt{║◌│◌│◌│◌║ 7} \\
\texttt{╚═╧═╧═╧═╝} \\
\texttt{```} \\
\texttt{Tá sibh ag imirt babhta 18 faoi láthair as uasmhéid 28.} \\
\texttt{Tá 14 as 16 pionós carntha agaibh le chéile.} \\
\texttt{Bhog an t{-}imreoir eile réad ar a ngreille.} \\
\\ 
\\ 
\texttt{Cad é do chéad ordú eile?} \\
            }
        }
    }
    & \\ \\

    \theutterance \stepcounter{utterance}  
    & \multicolumn{4}{p{0.6\linewidth}}{
        \cellcolor[rgb]{0.9,0.9,0.9}{
            \makecell[{{p{\linewidth}}}]{
                \texttt{\tiny{[P1$\rangle$GM]}}
                \texttt{BOG: M, (6,7)} \\
            }
        }
    }
    & & \\ \\

    \theutterance \stepcounter{utterance}  
    & & & \multicolumn{2}{p{0.3\linewidth}}{
        \cellcolor[rgb]{0.9,0.9,0.9}{
            \makecell[{{p{\linewidth}}}]{
                \texttt{\tiny{[GM$|$GM]}}
                \texttt{BOG: M, (6,7)} \\
            }
        }
    }
    & & \\ \\

    \theutterance \stepcounter{utterance}  
    & & & \multicolumn{2}{p{0.3\linewidth}}{
        \cellcolor[rgb]{0.9,0.9,0.9}{
            \makecell[{{p{\linewidth}}}]{
                \texttt{\tiny{[GM$|$GM]}}
                \texttt{Pionós: níl (6,7) folamh, ach tá '│' ann.} \\
\\ 
\texttt{Tá 15 as 16 pionós carntha agaibh le chéile.} \\
\texttt{Déan iarracht arís, le do thoil!} \\
            }
        }
    }
    & & \\ \\

    \theutterance \stepcounter{utterance}  
    & & \multicolumn{4}{p{0.6\linewidth}}{
        \cellcolor[rgb]{0.9,0.9,0.9}{
            \makecell[{{p{\linewidth}}}]{
                \texttt{\tiny{[P1$\langle$GM]}}
                \texttt{Pionós: níl (6,7) folamh, ach tá '│' ann.} \\
\\ 
\texttt{Tá 15 as 16 pionós carntha agaibh le chéile.} \\
\texttt{Déan iarracht arís, le do thoil!} \\
            }
        }
    }
    & \\ \\

    \theutterance \stepcounter{utterance}  
    & \multicolumn{4}{p{0.6\linewidth}}{
        \cellcolor[rgb]{0.9,0.9,0.9}{
            \makecell[{{p{\linewidth}}}]{
                \texttt{\tiny{[P1$\rangle$GM]}}
                \texttt{BOG: M, (7,7)} \\
            }
        }
    }
    & & \\ \\

    \theutterance \stepcounter{utterance}  
    & & & \multicolumn{2}{p{0.3\linewidth}}{
        \cellcolor[rgb]{0.9,0.9,0.9}{
            \makecell[{{p{\linewidth}}}]{
                \texttt{\tiny{[GM$|$GM]}}
                \texttt{BOG: M, (7,7)} \\
            }
        }
    }
    & & \\ \\

    \theutterance \stepcounter{utterance}  
    & & & \multicolumn{2}{p{0.3\linewidth}}{
        \cellcolor[rgb]{0.9,0.9,0.9}{
            \makecell[{{p{\linewidth}}}]{
                \texttt{\tiny{[GM$|$GM]}}
                \texttt{M moved to (7, 7)} \\
            }
        }
    }
    & & \\ \\

    \theutterance \stepcounter{utterance}  
    & & \multicolumn{4}{p{0.6\linewidth}}{
        \cellcolor[rgb]{0.9,0.9,0.9}{
            \makecell[{{p{\linewidth}}}]{
                \texttt{\tiny{[GM$\rangle$P2]}}
                \texttt{Bogadh 'M' go (5,7) go rathúil. Seo mar atá do ghreille nuashonraithe:} \\
\texttt{```} \\
\texttt{ 1234567} \\
\texttt{╔═╤══╤══╗} \\
\texttt{║◌│◌◌│◌◌║ 1} \\
\texttt{║◌└┬─┴──╢ 2} \\
\texttt{║A◌│◌OT◌║ 3} \\
\texttt{╟─┬┘SP┌─╢ 4} \\
\texttt{║◌│◌◌D│◌║ 5} \\
\texttt{╟─┼─┬─┤◌║ 6} \\
\texttt{║◌│◌│M│◌║ 7} \\
\texttt{╚═╧═╧═╧═╝} \\
\texttt{```} \\
\texttt{Tá sibh ag imirt babhta 19 faoi láthair as uasmhéid 28.} \\
\texttt{Tá 15 as 16 pionós carntha agaibh le chéile.} \\
\texttt{Bhog an t{-}imreoir eile réad ar a ngreille.} \\
\\ 
\\ 
\texttt{Cad é do chéad ordú eile?} \\
            }
        }
    }
    & \\ \\

    \theutterance \stepcounter{utterance}  
    & & & \multicolumn{4}{p{0.6\linewidth}}{
        \cellcolor[rgb]{0.9,0.9,0.9}{
            \makecell[{{p{\linewidth}}}]{
                \texttt{\tiny{[GM$\langle$P2]}}
                \texttt{ABAIR: D'éirigh liom M a bhogadh go (5,7) {-} ní raibh (6,7) folamh mar gheall ar '│'. Mar sin tá M ag (5,7) agam anois in ionad (6,7). An féidir leat do M a bhogadh go (5,7) ina ionad? Seo iad mo shuímh reatha: P(5,4), O(5,3), T(6,3), S(4,4), D(5,5), A(1,3), M(5,7). An bhfuilimid réidh le seiceáil an bhfuil gach rud ailínithe? Níl ach 1 phionós fágtha againn!} \\
            }
        }
    }
     \\ \\

    \theutterance \stepcounter{utterance}  
    & & & \multicolumn{2}{p{0.3\linewidth}}{
        \cellcolor[rgb]{0.9,0.9,0.9}{
            \makecell[{{p{\linewidth}}}]{
                \texttt{\tiny{[GM$|$GM]}}
                \texttt{ABAIR: D'éirigh liom M a bhogadh go (5,7) {-} ní raibh (6,7) folamh mar gheall ar '│'. Mar sin tá M ag (5,7) agam anois in ionad (6,7). An féidir leat do M a bhogadh go (5,7) ina ionad? Seo iad mo shuímh reatha: P(5,4), O(5,3), T(6,3), S(4,4), D(5,5), A(1,3), M(5,7). An bhfuilimid réidh le seiceáil an bhfuil gach rud ailínithe? Níl ach 1 phionós fágtha againn!} \\
            }
        }
    }
    & & \\ \\

    \theutterance \stepcounter{utterance}  
    & & \multicolumn{4}{p{0.6\linewidth}}{
        \cellcolor[rgb]{0.9,0.9,0.9}{
            \makecell[{{p{\linewidth}}}]{
                \texttt{\tiny{[P1$\langle$GM]}}
                \texttt{Bogadh 'M' go (7,7) go rathúil. Seo mar atá do ghreille nuashonraithe:} \\
\texttt{```} \\
\texttt{ 1234567} \\
\texttt{╔═╤══╤══╗} \\
\texttt{║◌│◌◌│◌◌║ 1} \\
\texttt{║T└┬─┴──╢ 2} \\
\texttt{║A◌│◌O◌◌║ 3} \\
\texttt{╟─┬┘SP┌─╢ 4} \\
\texttt{║◌│◌◌◌│D║ 5} \\
\texttt{╟─┼─┬─┤◌║ 6} \\
\texttt{║◌│◌│◌│M║ 7} \\
\texttt{╚═╧═╧═╧═╝} \\
\texttt{```} \\
\texttt{Tá sibh ag imirt babhta 19 faoi láthair as uasmhéid 28.} \\
\texttt{Tá 15 as 16 pionós carntha agaibh le chéile.} \\
\texttt{Sheol an t{-}imreoir eile an teachtaireacht seo:} \\
\\ 
\\ 
\texttt{"D'éirigh liom M a bhogadh go (5,7) {-} ní raibh (6,7) folamh mar gheall ar '│'. Mar sin tá M ag (5,7) agam anois in ionad (6,7). An féidir leat do M a bhogadh go (5,7) ina ionad? Seo iad mo shuímh reatha: P(5,4), O(5,3), T(6,3), S(4,4), D(5,5), A(1,3), M(5,7). An bhfuilimid réidh le seiceáil an bhfuil gach rud ailínithe? Níl ach 1 phionós fágtha againn!"} \\
\\ 
\\ 
\texttt{Cad é do chéad ordú eile?} \\
            }
        }
    }
    & \\ \\

    \theutterance \stepcounter{utterance}  
    & \multicolumn{4}{p{0.6\linewidth}}{
        \cellcolor[rgb]{0.9,0.9,0.9}{
            \makecell[{{p{\linewidth}}}]{
                \texttt{\tiny{[P1$\rangle$GM]}}
                \texttt{ABAIR: D'athraigh tú M go (5,7), ach an bhfuil (6,7) folamh ar do chlár? Ar mo chlár, tá (5,7) folamh. Faoi láthair tá mo M ag (7,7). Ní mór dúinn cinneadh a dhéanamh {-} an sprioc é M(5,7) nó M(6,7)? Tá (5,7) folamh agam, agus is féidir liom M a bhogadh ann. Ach tá 1 phionós amháin fágtha againn, mar sin caithfidh mé a bheith cinnte sula mbogfaidh mé. Deimhnigh: an é M(5,7) an sprioc deiridh anois? Agus an bhfuil gach rud eile ag na suímh chearta: P(5,4), O(5,3), T(6,3), S(4,4), D(5,5), A(1,3), M(5,7)?} \\
            }
        }
    }
    & & \\ \\

    \theutterance \stepcounter{utterance}  
    & & & \multicolumn{2}{p{0.3\linewidth}}{
        \cellcolor[rgb]{0.9,0.9,0.9}{
            \makecell[{{p{\linewidth}}}]{
                \texttt{\tiny{[GM$|$GM]}}
                \texttt{ABAIR: D'athraigh tú M go (5,7), ach an bhfuil (6,7) folamh ar do chlár? Ar mo chlár, tá (5,7) folamh. Faoi láthair tá mo M ag (7,7). Ní mór dúinn cinneadh a dhéanamh {-} an sprioc é M(5,7) nó M(6,7)? Tá (5,7) folamh agam, agus is féidir liom M a bhogadh ann. Ach tá 1 phionós amháin fágtha againn, mar sin caithfidh mé a bheith cinnte sula mbogfaidh mé. Deimhnigh: an é M(5,7) an sprioc deiridh anois? Agus an bhfuil gach rud eile ag na suímh chearta: P(5,4), O(5,3), T(6,3), S(4,4), D(5,5), A(1,3), M(5,7)?} \\
            }
        }
    }
    & & \\ \\

    \theutterance \stepcounter{utterance}  
    & & \multicolumn{4}{p{0.6\linewidth}}{
        \cellcolor[rgb]{0.9,0.9,0.9}{
            \makecell[{{p{\linewidth}}}]{
                \texttt{\tiny{[GM$\rangle$P2]}}
                \texttt{Cuireadh do theachtaireacht ar aghaidh chuig an imreoir eile.} \\
\\ 
\\ 
\texttt{Tá sibh ag imirt babhta 20 faoi láthair as uasmhéid 28.} \\
\texttt{Tá 15 as 16 pionós carntha agaibh le chéile.} \\
\texttt{Sheol an t{-}imreoir eile an teachtaireacht seo:} \\
\\ 
\\ 
\texttt{"D'athraigh tú M go (5,7), ach an bhfuil (6,7) folamh ar do chlár? Ar mo chlár, tá (5,7) folamh. Faoi láthair tá mo M ag (7,7). Ní mór dúinn cinneadh a dhéanamh {-} an sprioc é M(5,7) nó M(6,7)? Tá (5,7) folamh agam, agus is féidir liom M a bhogadh ann. Ach tá 1 phionós amháin fágtha againn, mar sin caithfidh mé a bheith cinnte sula mbogfaidh mé. Deimhnigh: an é M(5,7) an sprioc deiridh anois? Agus an bhfuil gach rud eile ag na suímh chearta: P(5,4), O(5,3), T(6,3), S(4,4), D(5,5), A(1,3), M(5,7)?"} \\
\\ 
\\ 
\texttt{Cad é do chéad ordú eile?} \\
            }
        }
    }
    & \\ \\

    \theutterance \stepcounter{utterance}  
    & & & \multicolumn{4}{p{0.6\linewidth}}{
        \cellcolor[rgb]{0.9,0.9,0.9}{
            \makecell[{{p{\linewidth}}}]{
                \texttt{\tiny{[GM$\langle$P2]}}
                \texttt{ABAIR: Deimhním: is é M(5,7) an sprioc deiridh. Tá gach rud eile ag na suímh chearta: P(5,4), O(5,3), T(6,3), S(4,4), D(5,5), A(1,3), M(5,7). Bog do M go (5,7) agus beidh muid réidh! Ná húsáid ach (5,7) {-} ná triail (6,7) mar níl sé folamh ar mo chlár. Tá 1 phionós fágtha againn, mar sin bí cúramach. Nuair a bheidh sé sin déanta, abair liom agus iarrfaidh mé críochnú!} \\
            }
        }
    }
     \\ \\

    \theutterance \stepcounter{utterance}  
    & & & \multicolumn{2}{p{0.3\linewidth}}{
        \cellcolor[rgb]{0.9,0.9,0.9}{
            \makecell[{{p{\linewidth}}}]{
                \texttt{\tiny{[GM$|$GM]}}
                \texttt{ABAIR: Deimhním: is é M(5,7) an sprioc deiridh. Tá gach rud eile ag na suímh chearta: P(5,4), O(5,3), T(6,3), S(4,4), D(5,5), A(1,3), M(5,7). Bog do M go (5,7) agus beidh muid réidh! Ná húsáid ach (5,7) {-} ná triail (6,7) mar níl sé folamh ar mo chlár. Tá 1 phionós fágtha againn, mar sin bí cúramach. Nuair a bheidh sé sin déanta, abair liom agus iarrfaidh mé críochnú!} \\
            }
        }
    }
    & & \\ \\

    \theutterance \stepcounter{utterance}  
    & & \multicolumn{4}{p{0.6\linewidth}}{
        \cellcolor[rgb]{0.9,0.9,0.9}{
            \makecell[{{p{\linewidth}}}]{
                \texttt{\tiny{[P1$\langle$GM]}}
                \texttt{Cuireadh do theachtaireacht ar aghaidh chuig an imreoir eile.} \\
\\ 
\\ 
\texttt{Tá sibh ag imirt babhta 20 faoi láthair as uasmhéid 28.} \\
\texttt{Tá 15 as 16 pionós carntha agaibh le chéile.} \\
\texttt{Sheol an t{-}imreoir eile an teachtaireacht seo:} \\
\\ 
\\ 
\texttt{"Deimhním: is é M(5,7) an sprioc deiridh. Tá gach rud eile ag na suímh chearta: P(5,4), O(5,3), T(6,3), S(4,4), D(5,5), A(1,3), M(5,7). Bog do M go (5,7) agus beidh muid réidh! Ná húsáid ach (5,7) {-} ná triail (6,7) mar níl sé folamh ar mo chlár. Tá 1 phionós fágtha againn, mar sin bí cúramach. Nuair a bheidh sé sin déanta, abair liom agus iarrfaidh mé críochnú!"} \\
\\ 
\\ 
\texttt{Cad é do chéad ordú eile?} \\
            }
        }
    }
    & \\ \\

    \theutterance \stepcounter{utterance}  
    & \multicolumn{4}{p{0.6\linewidth}}{
        \cellcolor[rgb]{0.9,0.9,0.9}{
            \makecell[{{p{\linewidth}}}]{
                \texttt{\tiny{[P1$\rangle$GM]}}
                \texttt{BOG: M, (5,7)} \\
            }
        }
    }
    & & \\ \\

    \theutterance \stepcounter{utterance}  
    & & & \multicolumn{2}{p{0.3\linewidth}}{
        \cellcolor[rgb]{0.9,0.9,0.9}{
            \makecell[{{p{\linewidth}}}]{
                \texttt{\tiny{[GM$|$GM]}}
                \texttt{BOG: M, (5,7)} \\
            }
        }
    }
    & & \\ \\

    \theutterance \stepcounter{utterance}  
    & & & \multicolumn{2}{p{0.3\linewidth}}{
        \cellcolor[rgb]{0.9,0.9,0.9}{
            \makecell[{{p{\linewidth}}}]{
                \texttt{\tiny{[GM$|$GM]}}
                \texttt{M moved to (5, 7)} \\
            }
        }
    }
    & & \\ \\

    \theutterance \stepcounter{utterance}  
    & & \multicolumn{4}{p{0.6\linewidth}}{
        \cellcolor[rgb]{0.9,0.9,0.9}{
            \makecell[{{p{\linewidth}}}]{
                \texttt{\tiny{[GM$\rangle$P2]}}
                \texttt{Cuireadh do theachtaireacht ar aghaidh chuig an imreoir eile.} \\
\\ 
\\ 
\texttt{Tá sibh ag imirt babhta 21 faoi láthair as uasmhéid 28.} \\
\texttt{Tá 15 as 16 pionós carntha agaibh le chéile.} \\
\texttt{Bhog an t{-}imreoir eile réad ar a ngreille.} \\
\\ 
\\ 
\texttt{Cad é do chéad ordú eile?} \\
            }
        }
    }
    & \\ \\

    \theutterance \stepcounter{utterance}  
    & & & \multicolumn{4}{p{0.6\linewidth}}{
        \cellcolor[rgb]{0.9,0.9,0.9}{
            \makecell[{{p{\linewidth}}}]{
                \texttt{\tiny{[GM$\langle$P2]}}
                \texttt{ABAIR: Ar bhog tú M go (5,7)? Má tá gach rud ag na suímh chearta anois, ba cheart dúinn an cluiche a chríochnú. An bhfuil tú réidh le críochnú? Deimhnigh le do thoil!} \\
            }
        }
    }
     \\ \\

    \theutterance \stepcounter{utterance}  
    & & & \multicolumn{2}{p{0.3\linewidth}}{
        \cellcolor[rgb]{0.9,0.9,0.9}{
            \makecell[{{p{\linewidth}}}]{
                \texttt{\tiny{[GM$|$GM]}}
                \texttt{ABAIR: Ar bhog tú M go (5,7)? Má tá gach rud ag na suímh chearta anois, ba cheart dúinn an cluiche a chríochnú. An bhfuil tú réidh le críochnú? Deimhnigh le do thoil!} \\
            }
        }
    }
    & & \\ \\

    \theutterance \stepcounter{utterance}  
    & & \multicolumn{4}{p{0.6\linewidth}}{
        \cellcolor[rgb]{0.9,0.9,0.9}{
            \makecell[{{p{\linewidth}}}]{
                \texttt{\tiny{[P1$\langle$GM]}}
                \texttt{Bogadh 'M' go (5,7) go rathúil. Seo mar atá do ghreille nuashonraithe:} \\
\texttt{```} \\
\texttt{ 1234567} \\
\texttt{╔═╤══╤══╗} \\
\texttt{║◌│◌◌│◌◌║ 1} \\
\texttt{║T└┬─┴──╢ 2} \\
\texttt{║A◌│◌O◌◌║ 3} \\
\texttt{╟─┬┘SP┌─╢ 4} \\
\texttt{║◌│◌◌◌│D║ 5} \\
\texttt{╟─┼─┬─┤◌║ 6} \\
\texttt{║◌│◌│M│◌║ 7} \\
\texttt{╚═╧═╧═╧═╝} \\
\texttt{```} \\
\texttt{Tá sibh ag imirt babhta 21 faoi láthair as uasmhéid 28.} \\
\texttt{Tá 15 as 16 pionós carntha agaibh le chéile.} \\
\texttt{Sheol an t{-}imreoir eile an teachtaireacht seo:} \\
\\ 
\\ 
\texttt{"Ar bhog tú M go (5,7)? Má tá gach rud ag na suímh chearta anois, ba cheart dúinn an cluiche a chríochnú. An bhfuil tú réidh le críochnú? Deimhnigh le do thoil!"} \\
\\ 
\\ 
\texttt{Cad é do chéad ordú eile?} \\
            }
        }
    }
    & \\ \\

    \theutterance \stepcounter{utterance}  
    & \multicolumn{4}{p{0.6\linewidth}}{
        \cellcolor[rgb]{0.9,0.9,0.9}{
            \makecell[{{p{\linewidth}}}]{
                \texttt{\tiny{[P1$\rangle$GM]}}
                \texttt{ABAIR: Níl mo rudaí go léir ag na suímh chearta fós. Tá T fós ag (1,2) agus D ag (7,5) agam. Caithfidh mé T a bhogadh go (6,3) agus D a bhogadh go (5,5). An bhfuil (6,3) agus (5,5) folamh ar do chlár faoi láthair? Tá T agat cheana ag (6,3) agus D ag (5,5), mar sin deimhnigh le do thoil go bhfanfaidh siad ann. Bogfaidh mé anois iad ceann ar cheann.} \\
            }
        }
    }
    & & \\ \\

    \theutterance \stepcounter{utterance}  
    & & & \multicolumn{2}{p{0.3\linewidth}}{
        \cellcolor[rgb]{0.9,0.9,0.9}{
            \makecell[{{p{\linewidth}}}]{
                \texttt{\tiny{[GM$|$GM]}}
                \texttt{ABAIR: Níl mo rudaí go léir ag na suímh chearta fós. Tá T fós ag (1,2) agus D ag (7,5) agam. Caithfidh mé T a bhogadh go (6,3) agus D a bhogadh go (5,5). An bhfuil (6,3) agus (5,5) folamh ar do chlár faoi láthair? Tá T agat cheana ag (6,3) agus D ag (5,5), mar sin deimhnigh le do thoil go bhfanfaidh siad ann. Bogfaidh mé anois iad ceann ar cheann.} \\
            }
        }
    }
    & & \\ \\

    \theutterance \stepcounter{utterance}  
    & & \multicolumn{4}{p{0.6\linewidth}}{
        \cellcolor[rgb]{0.9,0.9,0.9}{
            \makecell[{{p{\linewidth}}}]{
                \texttt{\tiny{[GM$\rangle$P2]}}
                \texttt{Cuireadh do theachtaireacht ar aghaidh chuig an imreoir eile.} \\
\\ 
\\ 
\texttt{Tá sibh ag imirt babhta 22 faoi láthair as uasmhéid 28.} \\
\texttt{Tá 15 as 16 pionós carntha agaibh le chéile.} \\
\texttt{Sheol an t{-}imreoir eile an teachtaireacht seo:} \\
\\ 
\\ 
\texttt{"Níl mo rudaí go léir ag na suímh chearta fós. Tá T fós ag (1,2) agus D ag (7,5) agam. Caithfidh mé T a bhogadh go (6,3) agus D a bhogadh go (5,5). An bhfuil (6,3) agus (5,5) folamh ar do chlár faoi láthair? Tá T agat cheana ag (6,3) agus D ag (5,5), mar sin deimhnigh le do thoil go bhfanfaidh siad ann. Bogfaidh mé anois iad ceann ar cheann."} \\
\\ 
\\ 
\texttt{Cad é do chéad ordú eile?} \\
            }
        }
    }
    & \\ \\

    \theutterance \stepcounter{utterance}  
    & & & \multicolumn{4}{p{0.6\linewidth}}{
        \cellcolor[rgb]{0.9,0.9,0.9}{
            \makecell[{{p{\linewidth}}}]{
                \texttt{\tiny{[GM$\langle$P2]}}
                \texttt{ABAIR: Deimhním go bhfuil T agam ag (6,3) agus D agam ag (5,5) agus fanfaidh siad ann. Tá siad sa phost cheana féin. Bog do T go (6,3) agus do D go (5,5) {-} beidh siad ag teacht le mo shuímh. Bí cúramach le do bhogtha {-} níl ach 1 phionós fágtha againn! Tar éis duit iad a bhogadh, beidh gach rud ailínithe: P(5,4), O(5,3), T(6,3), S(4,4), D(5,5), A(1,3), M(5,7). Abair liom nuair atá tú réidh!} \\
            }
        }
    }
     \\ \\

    \theutterance \stepcounter{utterance}  
    & & & \multicolumn{2}{p{0.3\linewidth}}{
        \cellcolor[rgb]{0.9,0.9,0.9}{
            \makecell[{{p{\linewidth}}}]{
                \texttt{\tiny{[GM$|$GM]}}
                \texttt{ABAIR: Deimhním go bhfuil T agam ag (6,3) agus D agam ag (5,5) agus fanfaidh siad ann. Tá siad sa phost cheana féin. Bog do T go (6,3) agus do D go (5,5) {-} beidh siad ag teacht le mo shuímh. Bí cúramach le do bhogtha {-} níl ach 1 phionós fágtha againn! Tar éis duit iad a bhogadh, beidh gach rud ailínithe: P(5,4), O(5,3), T(6,3), S(4,4), D(5,5), A(1,3), M(5,7). Abair liom nuair atá tú réidh!} \\
            }
        }
    }
    & & \\ \\

    \theutterance \stepcounter{utterance}  
    & & \multicolumn{4}{p{0.6\linewidth}}{
        \cellcolor[rgb]{0.9,0.9,0.9}{
            \makecell[{{p{\linewidth}}}]{
                \texttt{\tiny{[P1$\langle$GM]}}
                \texttt{Cuireadh do theachtaireacht ar aghaidh chuig an imreoir eile.} \\
\\ 
\\ 
\texttt{Tá sibh ag imirt babhta 22 faoi láthair as uasmhéid 28.} \\
\texttt{Tá 15 as 16 pionós carntha agaibh le chéile.} \\
\texttt{Sheol an t{-}imreoir eile an teachtaireacht seo:} \\
\\ 
\\ 
\texttt{"Deimhním go bhfuil T agam ag (6,3) agus D agam ag (5,5) agus fanfaidh siad ann. Tá siad sa phost cheana féin. Bog do T go (6,3) agus do D go (5,5) {-} beidh siad ag teacht le mo shuímh. Bí cúramach le do bhogtha {-} níl ach 1 phionós fágtha againn! Tar éis duit iad a bhogadh, beidh gach rud ailínithe: P(5,4), O(5,3), T(6,3), S(4,4), D(5,5), A(1,3), M(5,7). Abair liom nuair atá tú réidh!"} \\
\\ 
\\ 
\texttt{Cad é do chéad ordú eile?} \\
            }
        }
    }
    & \\ \\

    \theutterance \stepcounter{utterance}  
    & \multicolumn{4}{p{0.6\linewidth}}{
        \cellcolor[rgb]{0.9,0.9,0.9}{
            \makecell[{{p{\linewidth}}}]{
                \texttt{\tiny{[P1$\rangle$GM]}}
                \texttt{BOG: T, (6,3)} \\
            }
        }
    }
    & & \\ \\

    \theutterance \stepcounter{utterance}  
    & & & \multicolumn{2}{p{0.3\linewidth}}{
        \cellcolor[rgb]{0.9,0.9,0.9}{
            \makecell[{{p{\linewidth}}}]{
                \texttt{\tiny{[GM$|$GM]}}
                \texttt{BOG: T, (6,3)} \\
            }
        }
    }
    & & \\ \\

    \theutterance \stepcounter{utterance}  
    & & & \multicolumn{2}{p{0.3\linewidth}}{
        \cellcolor[rgb]{0.9,0.9,0.9}{
            \makecell[{{p{\linewidth}}}]{
                \texttt{\tiny{[GM$|$GM]}}
                \texttt{T moved to (6, 3)} \\
            }
        }
    }
    & & \\ \\

    \theutterance \stepcounter{utterance}  
    & & \multicolumn{4}{p{0.6\linewidth}}{
        \cellcolor[rgb]{0.9,0.9,0.9}{
            \makecell[{{p{\linewidth}}}]{
                \texttt{\tiny{[GM$\rangle$P2]}}
                \texttt{Cuireadh do theachtaireacht ar aghaidh chuig an imreoir eile.} \\
\\ 
\\ 
\texttt{Tá sibh ag imirt babhta 23 faoi láthair as uasmhéid 28.} \\
\texttt{Tá 15 as 16 pionós carntha agaibh le chéile.} \\
\texttt{Bhog an t{-}imreoir eile réad ar a ngreille.} \\
\\ 
\\ 
\texttt{Cad é do chéad ordú eile?} \\
            }
        }
    }
    & \\ \\

    \theutterance \stepcounter{utterance}  
    & & & \multicolumn{4}{p{0.6\linewidth}}{
        \cellcolor[rgb]{0.9,0.9,0.9}{
            \makecell[{{p{\linewidth}}}]{
                \texttt{\tiny{[GM$\langle$P2]}}
                \texttt{ABAIR: Fanfaidh mé go dtí go mbeidh do bhogtha déanta agat. Níl aon bhogadh le déanamh agam {-} tá mo rudaí go léir ag na sprioc{-}shuímh. Abair liom nuair atá T agus D bogtha agat, agus ansin iarrfaidh mé críochnú.} \\
            }
        }
    }
     \\ \\

    \theutterance \stepcounter{utterance}  
    & & & \multicolumn{2}{p{0.3\linewidth}}{
        \cellcolor[rgb]{0.9,0.9,0.9}{
            \makecell[{{p{\linewidth}}}]{
                \texttt{\tiny{[GM$|$GM]}}
                \texttt{ABAIR: Fanfaidh mé go dtí go mbeidh do bhogtha déanta agat. Níl aon bhogadh le déanamh agam {-} tá mo rudaí go léir ag na sprioc{-}shuímh. Abair liom nuair atá T agus D bogtha agat, agus ansin iarrfaidh mé críochnú.} \\
            }
        }
    }
    & & \\ \\

    \theutterance \stepcounter{utterance}  
    & & \multicolumn{4}{p{0.6\linewidth}}{
        \cellcolor[rgb]{0.9,0.9,0.9}{
            \makecell[{{p{\linewidth}}}]{
                \texttt{\tiny{[P1$\langle$GM]}}
                \texttt{Bogadh 'T' go (6,3) go rathúil. Seo mar atá do ghreille nuashonraithe:} \\
\texttt{```} \\
\texttt{ 1234567} \\
\texttt{╔═╤══╤══╗} \\
\texttt{║◌│◌◌│◌◌║ 1} \\
\texttt{║◌└┬─┴──╢ 2} \\
\texttt{║A◌│◌OT◌║ 3} \\
\texttt{╟─┬┘SP┌─╢ 4} \\
\texttt{║◌│◌◌◌│D║ 5} \\
\texttt{╟─┼─┬─┤◌║ 6} \\
\texttt{║◌│◌│M│◌║ 7} \\
\texttt{╚═╧═╧═╧═╝} \\
\texttt{```} \\
\texttt{Tá sibh ag imirt babhta 23 faoi láthair as uasmhéid 28.} \\
\texttt{Tá 15 as 16 pionós carntha agaibh le chéile.} \\
\texttt{Sheol an t{-}imreoir eile an teachtaireacht seo:} \\
\\ 
\\ 
\texttt{"Fanfaidh mé go dtí go mbeidh do bhogtha déanta agat. Níl aon bhogadh le déanamh agam {-} tá mo rudaí go léir ag na sprioc{-}shuímh. Abair liom nuair atá T agus D bogtha agat, agus ansin iarrfaidh mé críochnú."} \\
\\ 
\\ 
\texttt{Cad é do chéad ordú eile?} \\
            }
        }
    }
    & \\ \\

    \theutterance \stepcounter{utterance}  
    & \multicolumn{4}{p{0.6\linewidth}}{
        \cellcolor[rgb]{0.9,0.9,0.9}{
            \makecell[{{p{\linewidth}}}]{
                \texttt{\tiny{[P1$\rangle$GM]}}
                \texttt{BOG: D, (5,5)} \\
            }
        }
    }
    & & \\ \\

    \theutterance \stepcounter{utterance}  
    & & & \multicolumn{2}{p{0.3\linewidth}}{
        \cellcolor[rgb]{0.9,0.9,0.9}{
            \makecell[{{p{\linewidth}}}]{
                \texttt{\tiny{[GM$|$GM]}}
                \texttt{BOG: D, (5,5)} \\
            }
        }
    }
    & & \\ \\

    \theutterance \stepcounter{utterance}  
    & & & \multicolumn{2}{p{0.3\linewidth}}{
        \cellcolor[rgb]{0.9,0.9,0.9}{
            \makecell[{{p{\linewidth}}}]{
                \texttt{\tiny{[GM$|$GM]}}
                \texttt{D moved to (5, 5)} \\
            }
        }
    }
    & & \\ \\

    \theutterance \stepcounter{utterance}  
    & & \multicolumn{4}{p{0.6\linewidth}}{
        \cellcolor[rgb]{0.9,0.9,0.9}{
            \makecell[{{p{\linewidth}}}]{
                \texttt{\tiny{[GM$\rangle$P2]}}
                \texttt{Cuireadh do theachtaireacht ar aghaidh chuig an imreoir eile.} \\
\\ 
\\ 
\texttt{Tá sibh ag imirt babhta 24 faoi láthair as uasmhéid 28.} \\
\texttt{Tá 15 as 16 pionós carntha agaibh le chéile.} \\
\texttt{Bhog an t{-}imreoir eile réad ar a ngreille.} \\
\\ 
\\ 
\texttt{Cad é do chéad ordú eile?} \\
            }
        }
    }
    & \\ \\

    \theutterance \stepcounter{utterance}  
    & & & \multicolumn{4}{p{0.6\linewidth}}{
        \cellcolor[rgb]{0.9,0.9,0.9}{
            \makecell[{{p{\linewidth}}}]{
                \texttt{\tiny{[GM$\langle$P2]}}
                \texttt{ABAIR: Fanfaidh mé go foighneach. Níl aon bhogadh le déanamh agam {-} tá mo rudaí go léir sna suímh chomhaontaithe. Nuair a bheidh do chuid bogtha críochnaithe agat, abair liom agus déanfaimid críochnú. An bhfuil T agus D bogtha agat go fóill?} \\
            }
        }
    }
     \\ \\

    \theutterance \stepcounter{utterance}  
    & & & \multicolumn{2}{p{0.3\linewidth}}{
        \cellcolor[rgb]{0.9,0.9,0.9}{
            \makecell[{{p{\linewidth}}}]{
                \texttt{\tiny{[GM$|$GM]}}
                \texttt{ABAIR: Fanfaidh mé go foighneach. Níl aon bhogadh le déanamh agam {-} tá mo rudaí go léir sna suímh chomhaontaithe. Nuair a bheidh do chuid bogtha críochnaithe agat, abair liom agus déanfaimid críochnú. An bhfuil T agus D bogtha agat go fóill?} \\
            }
        }
    }
    & & \\ \\

    \theutterance \stepcounter{utterance}  
    & & \multicolumn{4}{p{0.6\linewidth}}{
        \cellcolor[rgb]{0.9,0.9,0.9}{
            \makecell[{{p{\linewidth}}}]{
                \texttt{\tiny{[P1$\langle$GM]}}
                \texttt{Bogadh 'D' go (5,5) go rathúil. Seo mar atá do ghreille nuashonraithe:} \\
\texttt{```} \\
\texttt{ 1234567} \\
\texttt{╔═╤══╤══╗} \\
\texttt{║◌│◌◌│◌◌║ 1} \\
\texttt{║◌└┬─┴──╢ 2} \\
\texttt{║A◌│◌OT◌║ 3} \\
\texttt{╟─┬┘SP┌─╢ 4} \\
\texttt{║◌│◌◌D│◌║ 5} \\
\texttt{╟─┼─┬─┤◌║ 6} \\
\texttt{║◌│◌│M│◌║ 7} \\
\texttt{╚═╧═╧═╧═╝} \\
\texttt{```} \\
\texttt{Tá sibh ag imirt babhta 24 faoi láthair as uasmhéid 28.} \\
\texttt{Tá 15 as 16 pionós carntha agaibh le chéile.} \\
\texttt{Sheol an t{-}imreoir eile an teachtaireacht seo:} \\
\\ 
\\ 
\texttt{"Fanfaidh mé go foighneach. Níl aon bhogadh le déanamh agam {-} tá mo rudaí go léir sna suímh chomhaontaithe. Nuair a bheidh do chuid bogtha críochnaithe agat, abair liom agus déanfaimid críochnú. An bhfuil T agus D bogtha agat go fóill?"} \\
\\ 
\\ 
\texttt{Cad é do chéad ordú eile?} \\
            }
        }
    }
    & \\ \\

    \theutterance \stepcounter{utterance}  
    & \multicolumn{4}{p{0.6\linewidth}}{
        \cellcolor[rgb]{0.9,0.9,0.9}{
            \makecell[{{p{\linewidth}}}]{
                \texttt{\tiny{[P1$\rangle$GM]}}
                \texttt{ABAIR: Tá gach rud bogtha agam anois! Seo iad mo shuímh: P(5,4), O(5,3), T(6,3), S(4,4), D(5,5), A(1,3), M(5,7). Ba cheart go mbeadh siad ag teacht le do shuímh. An bhfuilimid réidh le críochnú? críochnaithe?} \\
            }
        }
    }
    & & \\ \\

    \theutterance \stepcounter{utterance}  
    & & & \multicolumn{2}{p{0.3\linewidth}}{
        \cellcolor[rgb]{0.9,0.9,0.9}{
            \makecell[{{p{\linewidth}}}]{
                \texttt{\tiny{[GM$|$GM]}}
                \texttt{ABAIR: Tá gach rud bogtha agam anois! Seo iad mo shuímh: P(5,4), O(5,3), T(6,3), S(4,4), D(5,5), A(1,3), M(5,7). Ba cheart go mbeadh siad ag teacht le do shuímh. An bhfuilimid réidh le críochnú? críochnaithe?} \\
            }
        }
    }
    & & \\ \\

    \theutterance \stepcounter{utterance}  
    & & \multicolumn{4}{p{0.6\linewidth}}{
        \cellcolor[rgb]{0.9,0.9,0.9}{
            \makecell[{{p{\linewidth}}}]{
                \texttt{\tiny{[GM$\rangle$P2]}}
                \texttt{Cuireadh do theachtaireacht ar aghaidh chuig an imreoir eile.} \\
\\ 
\\ 
\texttt{Tá sibh ag imirt babhta 25 faoi láthair as uasmhéid 28.} \\
\texttt{Tá 15 as 16 pionós carntha agaibh le chéile.} \\
\texttt{Sheol an t{-}imreoir eile an teachtaireacht seo:} \\
\\ 
\\ 
\texttt{"Tá gach rud bogtha agam anois! Seo iad mo shuímh: P(5,4), O(5,3), T(6,3), S(4,4), D(5,5), A(1,3), M(5,7). Ba cheart go mbeadh siad ag teacht le do shuímh. An bhfuilimid réidh le críochnú? críochnaithe?"} \\
\\ 
\\ 
\texttt{Cad é do chéad ordú eile?} \\
            }
        }
    }
    & \\ \\

    \theutterance \stepcounter{utterance}  
    & & & \multicolumn{4}{p{0.6\linewidth}}{
        \cellcolor[rgb]{0.9,0.9,0.9}{
            \makecell[{{p{\linewidth}}}]{
                \texttt{\tiny{[GM$\langle$P2]}}
                \texttt{ABAIR: críochnaithe!} \\
            }
        }
    }
     \\ \\

    \theutterance \stepcounter{utterance}  
    & & & \multicolumn{2}{p{0.3\linewidth}}{
        \cellcolor[rgb]{0.9,0.9,0.9}{
            \makecell[{{p{\linewidth}}}]{
                \texttt{\tiny{[GM$|$GM]}}
                \texttt{ABAIR: críochnaithe!} \\
            }
        }
    }
    & & \\ \\

    \theutterance \stepcounter{utterance}  
    & & & \multicolumn{2}{p{0.3\linewidth}}{
        \cellcolor[rgb]{0.9,0.9,0.9}{
            \makecell[{{p{\linewidth}}}]{
                \texttt{\tiny{[GM$|$GM]}}
                \texttt{true} \\
            }
        }
    }
    & & \\ \\

    \theutterance \stepcounter{utterance}  
    & & & \multicolumn{2}{p{0.3\linewidth}}{
        \cellcolor[rgb]{0.9,0.9,0.9}{
            \makecell[{{p{\linewidth}}}]{
                \texttt{\tiny{[GM$|$GM]}}
                \texttt{Initial states:} \\
\texttt{```} \\
\texttt{Player 1:} \\
\texttt{ 1234567} \\
\texttt{╔═╤══╤══╗} \\
\texttt{║◌│◌T│◌◌║ 1} \\
\texttt{║◌└┬─┴──╢ 2} \\
\texttt{║◌◌│O◌◌◌║ 3} \\
\texttt{╟─┬┘P◌┌─╢ 4} \\
\texttt{║◌│AD◌│◌║ 5} \\
\texttt{╟─┼─┬─┤M║ 6} \\
\texttt{║◌│S│◌│◌║ 7} \\
\texttt{╚═╧═╧═╧═╝} \\
\\ 
\\ 
\texttt{Player 2:} \\
\texttt{ 1234567} \\
\texttt{╔═╤══╤══╗} \\
\texttt{║◌│O◌│◌◌║ 1} \\
\texttt{║◌└┬─┴──╢ 2} \\
\texttt{║◌P│◌D◌◌║ 3} \\
\texttt{╟─┬┘◌◌┌─╢ 4} \\
\texttt{║◌│◌◌◌│A║ 5} \\
\texttt{╟─┼─┬─┤S║ 6} \\
\texttt{║M│◌│◌│T║ 7} \\
\texttt{╚═╧═╧═╧═╝} \\
\\ 
\texttt{```} \\
            }
        }
    }
    & & \\ \\

    \theutterance \stepcounter{utterance}  
    & & & \multicolumn{2}{p{0.3\linewidth}}{
        \cellcolor[rgb]{0.9,0.9,0.9}{
            \makecell[{{p{\linewidth}}}]{
                \texttt{\tiny{[GM$|$GM]}}
                \texttt{End states:} \\
\texttt{```} \\
\texttt{Player 1:} \\
\texttt{ 1234567} \\
\texttt{╔═╤══╤══╗} \\
\texttt{║◌│◌◌│◌◌║ 1} \\
\texttt{║◌└┬─┴──╢ 2} \\
\texttt{║A◌│◌OT◌║ 3} \\
\texttt{╟─┬┘SP┌─╢ 4} \\
\texttt{║◌│◌◌D│◌║ 5} \\
\texttt{╟─┼─┬─┤◌║ 6} \\
\texttt{║◌│◌│M│◌║ 7} \\
\texttt{╚═╧═╧═╧═╝} \\
\\ 
\\ 
\texttt{Player 2:} \\
\texttt{ 1234567} \\
\texttt{╔═╤══╤══╗} \\
\texttt{║◌│◌◌│◌◌║ 1} \\
\texttt{║◌└┬─┴──╢ 2} \\
\texttt{║A◌│◌OT◌║ 3} \\
\texttt{╟─┬┘SP┌─╢ 4} \\
\texttt{║◌│◌◌D│◌║ 5} \\
\texttt{╟─┼─┬─┤◌║ 6} \\
\texttt{║◌│◌│M│◌║ 7} \\
\texttt{╚═╧═╧═╧═╝} \\
\\ 
\texttt{```} \\
            }
        }
    }
    & & \\ \\

    \theutterance \stepcounter{utterance}  
    & & & \multicolumn{2}{p{0.3\linewidth}}{
        \cellcolor[rgb]{0.9,0.9,0.9}{
            \makecell[{{p{\linewidth}}}]{
                \texttt{\tiny{[GM$|$GM]}}
                \texttt{* success: True} \\
\texttt{* lose: False} \\
\texttt{* aborted: False} \\
\texttt{{-}{-}{-}{-}{-}{-}{-}} \\
\texttt{* Shifts: 11.00} \\
\texttt{* Max Shifts: 12.00} \\
\texttt{* Min Shifts: 6.00} \\
\texttt{* End Distance Sum: 0.00} \\
\texttt{* Init Distance Sum: 27.62} \\
\texttt{* Expected Distance Sum: 29.33} \\
\texttt{* Penalties: 15.00} \\
\texttt{* Max Penalties: 16.00} \\
\texttt{* Rounds: 25.00} \\
\texttt{* Max Rounds: 28.00} \\
\texttt{* Object Count: 7.00} \\
            }
        }
    }
    & & \\ \\

\end{supertabular}
}

\end{document}
